% Sección incluida desde ARTURITO MAIN.tex con \documentclass[11pt,a4paper]{article}

\usepackage[utf8]{inputenc}
\usepackage[T1]{fontenc}
\usepackage[spanish,es-noquoting]{babel}
\usepackage[a4paper,margin=2.3cm]{geometry}
\usepackage{array}
\usepackage{longtable}
\usepackage{enumitem}
\usepackage{xcolor}
\usepackage{listings}
\usepackage{hyperref}
\usepackage{parskip}

\hypersetup{
  colorlinks=true,
  linkcolor=black,
  urlcolor=blue!45!black,
  pdftitle={Arturito --- Explicacion del firmware},
  pdfauthor={Equipo Arturito}
}

\urlstyle{tt}
\setlist{itemsep=2pt,topsep=4pt}
\lstset{
  basicstyle=\ttfamily\small,
  columns=fullflexible,
  keepspaces=true,
  breaklines=true,
  frame=single,
  rulecolor=\color{black!25}
}

\newcommand{\ruta}[1]{\path{#1}}
\newcommand{\capa}[1]{\textbf{Capa #1}}

\title{Arturito --- Explicacion del firmware\\
       \large Guia conceptual del workspace, los modulos y el flujo de datos}
\author{Equipo Arturito}
\date{3 de octubre de 2026}

\begin{document}
\maketitle

\begin{abstract}
Este documento explica como esta organizado y como funciona el firmware del
rover Arturito. Esta escrito para alguien que empieza desde cero: presenta el
workspace, las tres capas de software, las diferencias entre ESP32 y EDU-CIAA,
el runtime FreeRTOS compartido, el recorrido de una muestra desde los sensores
hasta Processing y las formas de compilar cada variante.

La explicacion describe el codigo actual. ESP32 ya fue probado con Processing.
La version EDU-CIAA compila y enlaza, pero todavia necesita una prueba fisica de
sensores en la placa. El diseno academico original se conserva por separado;
las decisiones vigentes se registran en \ruta{informe_avance.tex}.
\end{abstract}

\tableofcontents
\newpage

%==============================================================================
\section{Como leer esta guia}
%==============================================================================

El firmware puede entenderse como una cadena de responsabilidades:

\begin{center}
Sensor fisico $\rightarrow$ HAL de la placa $\rightarrow$ driver compartido
$\rightarrow$ tareas de aplicacion $\rightarrow$ HAL de telemetria
$\rightarrow$ USB/UART $\rightarrow$ Processing
\end{center}

Cada flecha representa una interfaz. El modulo de la izquierda ofrece una
funcion; el de la derecha la llama sin necesitar conocer sus detalles internos.
La meta es que los algoritmos que no dependen del hardware se escriban una sola
vez y que cada placa implemente unicamente los accesos fisicos que necesita.

Las palabras ``compartido'' y ``comun'' significan que el codigo se usa al
compilar para ambas placas. No significan que los perifericos internos de ESP32
y EDU-CIAA sean iguales. El MPU6050, el VL53L0X, las unidades y el protocolo de
la aplicacion son los mismos; las bibliotecas de bajo nivel y el arranque no.

%==============================================================================
\section{Que hace el firmware hoy}
%==============================================================================

La version actual realiza un ciclo de percepcion V0:

\begin{enumerate}
  \item Inicializa la IMU MPU6050 y el sensor de distancia VL53L0X.
  \item Lee la velocidad angular del eje Z de la IMU y una distancia del ToF.
  \item Corrige e integra la velocidad angular para estimar el angulo $\theta$.
  \item Acepta distancias entre 0.05 m y 1.50 m; el resto se representa como
        $-1$ (medida invalida o sensor no disponible).
  \item Publica una linea CSV con $\theta$ en radianes y distancia en metros.
  \item Processing representa la orientacion y acumula puntos del ToF mientras
        el robot gira en el lugar.
\end{enumerate}

No hay todavia odometria XY, encoders, motores, navegacion ni grilla de
ocupacion. El angulo proviene de integrar el giroscopio, por lo que deriva con
el tiempo. Esa limitacion es conocida y no equivale a un sistema SLAM.

%==============================================================================
\section{Mapa del workspace}
%==============================================================================

La raiz del repositorio es \ruta{C:/Arturito}. El archivo
\ruta{Arturito.code-workspace} abre cinco carpetas en VS Code:

\begin{lstlisting}
Arturito/
  Arturito.code-workspace
  README.md
  docs/
    ARTURITO MAIN.tex
    CONTEXTO_PROYECTO_ROVER.tex
    % (Diseno preliminar de firmware.tex)
    informe_avance.tex
    explicando_firmware.tex
  shared/
        config/config.h
    layer1_hal/
                  hal_imu.h
                  hal_time.h
                  hal_tof.h
                  hal_telemetry.h
    layer2_drivers/
                  driver_mpu6050.c/.h
                  driver_distancia.c/.h
    layer3_app/
                  app_runtime.c/.h
  esp32/
    platformio.ini
    src/
      main.cpp
      layer1_hal/
                        hal_imu_esp32.cpp
                        hal_time_esp32.cpp
                        hal_tof_esp32.cpp
                        hal_telemetry_esp32.cpp
      tests/
                        imu_only.cpp
                        tof_only.cpp
  edu-ciaa/
            Makefile, board.mk y program.mk
    rover/
      config.mk
      inc/
                        FreeRTOSConfig.h
      src/
        main.c
        layer1_hal/
                              hal_imu_educiaa.cpp
                              hal_time_educiaa.cpp
                              hal_tof_educiaa.cpp
                              hal_telemetry_educiaa.cpp
  processing/
    flecha/
    flecha_y_mapa_v1/
\end{lstlisting}

\begin{description}
  \item[\ruta{docs/}] Documentos del equipo. \ruta{informe_avance.tex} es la
        referencia viva de decisiones y estado; MAIN se conserva como entrega
        academica congelada.
  \item[\ruta{shared/}] Contratos HAL, drivers y logica de aplicacion que no
        deben depender de Arduino ni de sAPI.
  \item[\ruta{esp32/}] Firmware PlatformIO para ESP32, sus HAL de placa y
        programas de diagnostico independientes.
  \item[\ruta{edu-ciaa/}] Proyecto Make/LPCOpen/sAPI para EDU-CIAA. Adapta la
        misma aplicacion y drivers a su arquitectura.
  \item[\ruta{processing/}] Sketches que reciben el CSV por puerto serie y lo
        visualizan.
\end{description}

Las carpetas \ruta{layer3\_robotics/} y \ruta{layer4\_app/} pertenecian a una
organizacion anterior y fueron retiradas. La arquitectura acordada tiene tres
capas; su carpeta de aplicacion es \ruta{shared/layer3\_app/}.

%==============================================================================
\section{Las tres capas}
%==============================================================================

\subsection{Capa 1: HAL de perifericos}

HAL significa \emph{Hardware Abstraction Layer}, o capa de abstraccion de
hardware. Convierte operaciones de una placa concreta en funciones estables
que el resto del programa puede usar.

Los contratos publicos viven en \ruta{shared/layer1\_hal/}. Cada placa tiene
una implementacion con el mismo nombre de funcion:

\begin{itemize}
  \item \texttt{hal\_imu\_init()} inicializa la IMU y devuelve si pudo hacerlo.
  \item \texttt{hal\_imu\_get\_gyro\_z\_rads()} devuelve la velocidad angular
        Z en radianes por segundo.
  \item \texttt{hal\_tof\_init()} inicializa el sensor ToF.
  \item \texttt{hal\_tof\_get\_distance\_m()} devuelve metros, o $-1$ si la
        medida no esta disponible.
  \item \texttt{hal\_millis()} y \texttt{hal\_delay()} ofrecen tiempo y espera
        sin que los drivers llamen directamente a la API de una placa.
  \item \texttt{hal\_telemetry\_write()} envia una linea de texto por el canal
        serie de la placa.
\end{itemize}

La HAL no calcula el angulo del robot, no filtra la distancia y no arma el
mapa. Su responsabilidad es inicializar y leer el hardware, devolviendo
magnitudes con unidades acordadas.

\subsection{Capa 2: drivers}

Los drivers estan en \ruta{shared/layer2\_drivers/}. Conocen el significado
fisico de las muestras, pero usan solamente los contratos HAL.

\begin{itemize}
  \item \texttt{driver\_mpu6050.c} calibra el offset del giroscopio, aplica una
        zona muerta y convierte las muestras en un angulo integrado.
  \item \texttt{driver\_distancia.c} valida que una distancia este dentro del
        rango de trabajo del proyecto.
\end{itemize}

Al estar escritos contra las interfaces compartidas, estos drivers se compilan
sin duplicarlos para ESP32 y EDU-CIAA. La HAL correspondiente queda enlazada
por el build de cada placa.

\subsection{Capa 3: aplicacion compartida}

\ruta{shared/layer3\_app/app\_runtime.c} implementa el runtime FreeRTOS comun:
la inicializacion del sistema y las tareas de adquisicion, actualizacion de
drivers y publicacion. No accede a registros I2C ni llama directamente a
Arduino o sAPI.

La capa 3 todavia no implementa cinematica diferencial, pose XY, grilla ni
navegacion. Esas funciones dependen de encoders y decisiones de V1.

%==============================================================================
\section{Contratos compartidos}
%==============================================================================

\subsection{Configuracion comun}

\ruta{shared/config/config.h} define valores que deben coincidir en ambos
builds:

\begin{itemize}
  \item \texttt{SAMPLE\_PERIOD\_MS = 30}: periodo nominal de las tareas ciclicas.
  \item \texttt{IMU\_I2C\_ADDR = 0x68}: direccion I2C del MPU6050.
  \item \texttt{DIST\_MIN\_M = 0.05} y \texttt{DIST\_MAX\_M = 1.50}: limites
        aceptados para distancia.
\end{itemize}

El build define la placa con \texttt{BOARD\_ESP32} o
\texttt{BOARD\_EDUCIAA}. El header comun no elige hardware por su cuenta.

\subsection{Por que los headers tienen extern C}

Los contratos HAL y drivers estan escritos para poder llamarse desde C y C++.
Los headers rodean sus declaraciones con \texttt{extern "C"} cuando los incluye
C++, evitando que el compilador cambie los nombres de las funciones (\emph{name
mangling}). Asi, el codigo C compartido puede enlazarse con las implementaciones
HAL escritas en C++.

\subsection{Header agrupador}

\ruta{shared/layer1\_hal/hal\_sensors.h} incluye los headers separados de tiempo,
IMU y ToF. Es un header agrupador de compatibilidad; las funciones nuevas
incluyen directamente el contrato especifico que necesitan.

%==============================================================================
\section{Que hace cada modulo}
%==============================================================================

\subsection{Runtime: app\_runtime.c}

\begin{description}
  \item[\texttt{app\_runtime\_start()}] Inicializa el estado en valores seguros
        y crea la tarea bootstrap. Devuelve si FreeRTOS pudo crearla; no ejecuta
        la inicializacion de sensores directamente.
  \item[\texttt{task\_start\_application()}] Publica un encabezado de protocolo,
        inicializa IMU y ToF por separado, calibra el MPU6050, crea mutexes y
        luego crea las tres tareas ciclicas. Si la IMU falla, se detiene el
        runtime de sensores; si falla el ToF, se continua con distancia $-1$.
  \item[\texttt{task\_acquire()}] Lee gyro Z y distancia cruda por HAL y guarda
        una instantanea protegida por mutex. Espera el periodo nominal.
  \item[\texttt{task\_update\_drivers()}] Copia la ultima muestra cruda, integra
        gyro para obtener $\theta$, valida la distancia y guarda el estado de
        aplicacion bajo un segundo mutex.
  \item[\texttt{task\_publish()}] Copia el estado, forma
        \texttt{theta,distancia} con cuatro decimales y lo envia mediante la HAL
        de telemetria. Los retardos son nominales; no constituyen sincronizacion
        exacta de fase.
  \item[\texttt{stack\_depth\_from\_bytes()}] Recibe tamanos conceptuales en
        bytes. En ESP32 conserva ese valor, porque su port de FreeRTOS usa bytes;
        en CIAA lo convierte a palabras \texttt{StackType\_t}, como requiere el
        port estandar ARM.
  \item[\texttt{delete\_task\_if\_created()}] Libera tareas ya creadas si la
        creacion de una tarea posterior falla.
\end{description}

Los mutexes protegen las copias compartidas: el lector no observa una estructura
mientras otra tarea la esta reemplazando. No protegen los buses I2C como un
recurso compartido general; la aplicacion actual asigna las lecturas al flujo
de adquisicion.

\subsection{Driver MPU6050}

\begin{description}
  \item[\texttt{mpu6050\_init()}] Actualmente no hace trabajo: la
        inicializacion fisica ya la hizo \texttt{hal\_imu\_init()}.
  \item[\texttt{mpu6050\_calibrate()}] Espera que el robot quede quieto,
        promedia 1000 lecturas separadas por 2 ms para estimar el offset, y
        reinicia el estado de integracion.
  \item[\texttt{mpu6050\_update\_gyro(gz)}] Resta el offset, elimina velocidades
        dentro de $\pm1^\circ$/s, calcula un promedio trapezoidal entre muestras
        y acumula el angulo con el tiempo transcurrido.
  \item[\texttt{mpu6050\_update()}] Conveniencia para un super-loop o test:
        lee la HAL y delega en \texttt{mpu6050\_update\_gyro()}.
  \item[\texttt{mpu6050\_get\_theta\_rad()}] Devuelve el angulo integrado en
        radianes. La integracion de gyro deriva; no es una pose absoluta.
\end{description}

\subsection{Driver de distancia}

\begin{description}
  \item[\texttt{distancia\_init()}] Actualmente no hace trabajo; la HAL
        inicializa el sensor fisico.
  \item[\texttt{distancia\_validate\_m(d)}] Conserva $d$ si esta entre los
        limites inclusivos; devuelve $-1$ si esta fuera de rango.
  \item[\texttt{distancia\_get\_m()}] Lee la HAL ToF y valida el resultado en
        una sola llamada.
\end{description}

\subsection{Arranques de cada plataforma}

\begin{description}
  \item[ESP32, \texttt{setup()}] Arduino inicia el runtime FreeRTOS antes de
        llamar a \texttt{setup()}. El punto de entrada configura el baud rate y
        solicita \texttt{app\_runtime\_start()}. \texttt{loop()} queda vacio
        porque las tareas ejecutan la aplicacion.
  \item[EDU-CIAA, \texttt{main()}] Ejecuta \texttt{boardConfig()}, crea la tarea
        bootstrap con \texttt{app\_runtime\_start()} y llama a
        \texttt{vTaskStartScheduler()}. La inicializacion y calibracion ocurren
        despues, desde la tarea bootstrap. Si el scheduler retorna, informa un
        error porque normalmente no deberia retornar.
\end{description}

%==============================================================================
\section{FreeRTOS en las dos placas}
%==============================================================================

Las dos placas ejecutan la misma aplicacion FreeRTOS. Las cuatro tareas tienen
responsabilidades separadas:

\begin{longtable}{>{\raggedright\arraybackslash}p{0.27\textwidth}
                  >{\raggedright\arraybackslash}p{0.63\textwidth}}
\textbf{Tarea} & \textbf{Responsabilidad} \\
\hline
\texttt{app\_start} & Inicializacion de sensores, calibracion y creacion de
mutexes/tareas. Se elimina cuando termina el arranque. \\
\texttt{adquirir} & Lee IMU y ToF por HAL y publica una muestra cruda compartida. \\
\texttt{drivers} & Procesa las muestras crudas: offset, zona muerta, integracion
de $\theta$ y validacion de distancia. \\
\texttt{publicar} & Serializa el estado como CSV y delega la escritura a la HAL
de telemetria. \\
\end{longtable}

Las tareas ciclicas usan periodos nominales de 30 ms. Las prioridades son
arranque 4, adquisicion 3, drivers 2 y publicacion 1. CIAA admite hasta siete
niveles en su configuracion actual. Los valores de stack se expresan en bytes en
el runtime y se convierten para el ABI de FreeRTOS de CIAA.

La tarea de adquisicion realiza lecturas I2C secuenciales. La aplicacion no
promete que todas las etapas tengan latencia identica en ambas placas: el ToF
EDU-CIAA usa medicion single-shot y puede ocupar mas tiempo que una lectura
instantanea. FreeRTOS permite que otras tareas progresen, pero no vuelve
instantaneo al periferico.

\subsection{Scheduler y arranque}

\begin{itemize}
  \item En ESP32, Arduino/ESP-IDF arranca FreeRTOS. Crear tareas desde
        \texttt{setup()} no inicia un segundo scheduler.
  \item En EDU-CIAA, \texttt{main()} prepara el hardware, crea el bootstrap y
        llama a \texttt{vTaskStartScheduler()}.
  \item El bootstrap corre con el scheduler activo en ambas placas. Esto es
        necesario porque la calibracion usa \texttt{hal\_delay()}.
  \item CIAA compila el kernel FreeRTOS del repositorio, activa asignacion
        dinamica y utiliza \texttt{heap\_4}. En el estado actual, el heap se
        configura en 24 KiB, las prioridades disponibles son 7 y el tick es de
        1 kHz. La configuracion debe respetar la memoria disponible de LPC4337.
\end{itemize}

Los entry points son distintos porque Arduino y el proyecto LPCOpen/sAPI tienen
ciclos de inicio diferentes. Lo que se comparte es el runtime de aplicacion,
las responsabilidades de las tareas y los contratos; no se fuerza a ambas
plataformas a tener el mismo \texttt{main()} de sistema.

%==============================================================================
\section{HAL de ESP32}
%==============================================================================

\subsection{IMU: hal\_imu\_esp32.cpp}

Usa \texttt{Wire} para I2C. \texttt{hal\_imu\_init()} activa el MPU6050,
consulta su registro de identidad y falla si no recibe una respuesta valida.
\texttt{hal\_imu\_get\_gyro\_z\_rads()} lee los registros de gyro Z, aplica
la sensibilidad de $131$ cuentas por grado/segundo para el rango configurado y
convierte a radianes por segundo.

\subsection{Tiempo: hal\_time\_esp32.cpp}

Adapta \texttt{hal\_delay()} a \texttt{delay()} de Arduino y
\texttt{hal\_millis()} a \texttt{millis()}. El driver MPU usa estas funciones
para calibrar y calcular el intervalo $dt$.

\subsection{ToF: hal\_tof\_esp32.cpp}

Usa la biblioteca Adafruit \texttt{VL53L0X}. \texttt{hal\_tof\_init()} guarda
si la biblioteca encontro e inicializo el sensor. La lectura devuelve metros;
si el dispositivo no inicio o el estado de rango reporta error, devuelve $-1$.

\subsection{Telemetria: hal\_telemetry\_esp32.cpp}

Escribe cada linea con \texttt{Serial.println()}. No genera el contenido del
CSV: recibe texto ya formado desde la aplicacion.

%==============================================================================
\section{HAL de EDU-CIAA}
%==============================================================================

\subsection{IMU y tiempo}

\texttt{hal\_imu\_educiaa.cpp} utiliza las funciones MPU60X0 de sAPI para
inicializar, actualizar y leer gyro Z. \texttt{hal\_time\_educiaa.cpp} adapta
\texttt{hal\_delay()} a la espera sAPI y \texttt{hal\_millis()} a
\texttt{tickRead()}. Con el tick configurado a 1 kHz, el conteo representa
milisegundos.

\subsection{ToF EDU-CIAA}

\texttt{hal\_tof\_educiaa.cpp} implementa el acceso al VL53L0X mediante I2C de
sAPI. A diferencia de ESP32, no usa la biblioteca Adafruit: reproduce la
secuencia de registros del driver de referencia para la medicion single-shot.

Los auxiliares internos leen/escriben registros de 8 y 16 bits, traducen
unidades de timeout, consultan la secuencia de medicion y ejecutan la
configuracion SPAD/tuning/calibracion de referencia. La inicializacion comprueba
identificacion de modelo/revision y limita esperas mediante timeout. La lectura
inicia una medicion, espera su terminacion, lee milimetros y devuelve metros.
Ante fallo de identificacion, calibracion o timeout, marca el ToF no disponible
para que la aplicacion publique $-1$ y la IMU continue.

Esta implementacion compila y enlaza para EDU-CIAA. Sin sensor y placa conectados
no se puede asegurar que la secuencia I2C funcione en el montaje real.

\subsection{Telemetria: hal\_telemetry\_educiaa.cpp}

Implementa el mismo contrato \texttt{hal\_telemetry\_write()} con
\texttt{printf()} de sAPI y terminador CRLF. Asi el runtime compartido no sabe
si debajo escribe Arduino Serial o la consola serie de CIAA.

%==============================================================================
\section{Builds y perfiles}
%==============================================================================

\subsection{ESP32 / PlatformIO}

\ruta{esp32/platformio.ini} establece \texttt{esp32dev} como entorno por
defecto:

\begin{lstlisting}
cd esp32
pio run
pio run -t upload
\end{lstlisting}

Los perfiles \texttt{imu\_only} y \texttt{tof\_only} son diagnosticos
independientes; se seleccionan explicitamente con
\texttt{pio run -e imu\_only} o \texttt{pio run -e tof\_only}. Un comando de
upload con un perfil reemplaza el programa anterior en la placa: no se mezclan
los tres binarios.

\subsection{EDU-CIAA / Make}

\ruta{edu-ciaa/rover/config.mk} selecciona HAL EDU-CIAA, drivers compartidos y
FreeRTOS. La opcion \texttt{USE\_FREERTOS=y} incorpora el kernel; el tipo de
heap selecciona \texttt{heap\_4}. \ruta{edu-ciaa/rover/inc/FreeRTOSConfig.h}
define prioridades, frecuencia de tick, heap y opciones de chequeo.

\begin{lstlisting}
make -C edu-ciaa clean all
\end{lstlisting}

Este comando compila y enlaza; no carga firmware. El target \texttt{clean}
elimina productos de compilacion de \ruta{edu-ciaa/rover/out/}. No borra los
fuentes.

\subsection{Diagnosticos}

\begin{description}
  \item[\texttt{imu\_only}] Compila IMU + tiempo + driver MPU. Publica
        \texttt{theta\_rad,-1} sin depender del ToF ni del runtime completo.
  \item[\texttt{tof\_only}] Compila ToF + driver de distancia. Publica
        \texttt{0,distancia\_m} para comprobar el sensor de forma aislada.
\end{description}

Estos programas validan perifericos concretos, no prueban la planificacion de
las tres tareas del producto.

%==============================================================================
\section{Protocolo serie y Processing}
%==============================================================================

La aplicacion publica a 115200 baudios. Las lineas informativas comienzan con
\texttt{\#}; los sketch las ignoran. Cada muestra de datos tiene exactamente
dos columnas:

\begin{lstlisting}
theta_rad,distancia_m
\end{lstlisting}

$\theta$ esta en radianes; distancia esta en metros. El valor $-1$ indica que
no hay una medida ToF valida. No hay timestamp, posicion XY ni una tercera
columna.

\ruta{processing/flecha/flecha.pde} muestra la orientacion usando theta.
\ruta{processing/flecha\_y\_mapa\_v1/flecha\_y\_mapa\_v1.pde} lee las dos
columnas, descarta logs y datos no numericos, acepta distancias de 0.05 a 1.50 m
y las escala para dibujar una nube polar de hasta 400 puntos. Asume que el rover
gira en torno al origen; no es una grilla ni un mapa en coordenadas mundo.

Los sketches usan \texttt{COM3} como puerto por defecto. Cada PC debe usar el
puerto que realmente asigna al ESP32. Un puerto serie solo puede estar abierto
por una aplicacion a la vez; cerrar el monitor serie de PlatformIO antes de
abrir Processing.

%==============================================================================
\section{Recorrido de una muestra}
%==============================================================================

\begin{enumerate}
  \item El framework de la placa entra en su arranque: Arduino llama a
        \texttt{setup()}, o CIAA ejecuta \texttt{main()}.
  \item El adaptador de placa configura consola/placa y llama a
        \texttt{app\_runtime\_start()}.
  \item CIAA arranca el scheduler. El bootstrap comun inicializa IMU y ToF,
        calibra el MPU y crea los dos mutexes.
  \item Se crea \texttt{adquirir}. Lee $\omega_z$ y distancia por contratos HAL
        y actualiza \texttt{raw\_sample}.
  \item \texttt{drivers} toma una copia protegida, actualiza $\theta$ con la
        muestra IMU y valida distancia. Guarda \texttt{app\_sample}.
  \item \texttt{publicar} copia ese estado, forma el CSV y llama a
        \texttt{hal\_telemetry\_write()}.
  \item La implementacion serie de la placa escribe la linea; Processing la
        interpreta y dibuja la orientacion/medicion.
\end{enumerate}

Si falla la IMU, el bootstrap informa error y no crea las tareas de datos. Si
falla solo el ToF, la tarea de adquisicion sigue leyendo la IMU y el driver de
distancia transforma las medidas no disponibles en $-1$.

%==============================================================================
\section{Estado de validacion y limites}
%==============================================================================

\begin{itemize}
  \item Los entornos PlatformIO \texttt{esp32dev}, \texttt{imu\_only} y
        \texttt{tof\_only} compilan.
  \item El firmware ESP32 fue cargado y el plot V0 habia sido validado con
        Processing antes de compartir el runtime FreeRTOS.
  \item El build limpio EDU-CIAA con FreeRTOS compila y enlaza. No se cargo ni
        se probo con sensores en la placa EDU-CIAA.
  \item La carga del runtime compartido en ESP32 requiere confirmar telemetria
        en serie despues de la ultima modificacion; no se considera prueba de
        ejecucion solo porque haya compilado.
\end{itemize}

Limites funcionales actuales:

\begin{itemize}
  \item El angulo deriva porque se integra un giroscopio; no hay fusion con
        acelerometro ni correccion absoluta.
  \item Sin AS5600 no se calcula desplazamiento XY. El robot se supone en el
        origen para el plot polar V0.
  \item No hay motores, navegacion, grilla V1 ni SLAM implementados.
  \item La inicializacion ToF y los tiempos de medicion de las dos HAL son
        diferentes y requieren validacion fisica en cada placa.
\end{itemize}

%==============================================================================
\section{Glosario}
%==============================================================================

\begin{description}
  \item[HAL] Capa que traduce una API estable a operaciones de hardware de una
        placa concreta.
  \item[Driver] Codigo que transforma muestras HAL en magnitudes utiles para
        la aplicacion, aplicando calibracion, filtrado o validacion.
  \item[FreeRTOS] Sistema operativo en tiempo real que planifica las tareas.
  \item[Scheduler] Planificador que decide que tarea ejecuta el procesador.
  \item[Tarea] Funcion que FreeRTOS ejecuta como flujo independiente, normalmente
        dentro de un bucle.
  \item[Mutex] Cerradura que coordina acceso a datos compartidos entre tareas.
  \item[ToF] Sensor de distancia por tiempo de vuelo; aqui, VL53L0X.
  \item[CSV] Datos separados por comas; el protocolo V0 tiene dos columnas.
  \item[Stack] Memoria temporal de cada tarea para llamadas y variables locales.
\end{description}

% Sección incluida desde ARTURITO MAIN.tex con \documentclass[11pt,a4paper]{article}

\usepackage[utf8]{inputenc}
\usepackage[T1]{fontenc}
\usepackage[spanish,es-noquoting]{babel}
\usepackage[a4paper,margin=2.3cm]{geometry}
\usepackage{array}
\usepackage{longtable}
\usepackage{enumitem}
\usepackage{xcolor}
\usepackage{listings}
\usepackage{hyperref}
\usepackage{parskip}

\hypersetup{
  colorlinks=true,
  linkcolor=black,
  urlcolor=blue!45!black,
  pdftitle={Arturito --- Explicacion del firmware},
  pdfauthor={Equipo Arturito}
}

\urlstyle{tt}
\setlist{itemsep=2pt,topsep=4pt}
\lstset{
  basicstyle=\ttfamily\small,
  columns=fullflexible,
  keepspaces=true,
  breaklines=true,
  frame=single,
  rulecolor=\color{black!25}
}

\newcommand{\ruta}[1]{\path{#1}}
\newcommand{\capa}[1]{\textbf{Capa #1}}

\title{Arturito --- Explicacion del firmware\\
       \large Guia conceptual del workspace, los modulos y el flujo de datos}
\author{Equipo Arturito}
\date{3 de octubre de 2026}

\begin{document}
\maketitle

\begin{abstract}
Este documento explica como esta organizado y como funciona el firmware del
rover Arturito. Esta escrito para alguien que empieza desde cero: presenta el
workspace, las tres capas de software, las diferencias entre ESP32 y EDU-CIAA,
el runtime FreeRTOS compartido, el recorrido de una muestra desde los sensores
hasta Processing y las formas de compilar cada variante.

La explicacion describe el codigo actual. ESP32 ya fue probado con Processing.
La version EDU-CIAA compila y enlaza, pero todavia necesita una prueba fisica de
sensores en la placa. El diseno academico original se conserva por separado;
las decisiones vigentes se registran en \ruta{informe_avance.tex}.
\end{abstract}

\tableofcontents
\newpage

%==============================================================================
\section{Como leer esta guia}
%==============================================================================

El firmware puede entenderse como una cadena de responsabilidades:

\begin{center}
Sensor fisico $\rightarrow$ HAL de la placa $\rightarrow$ driver compartido
$\rightarrow$ tareas de aplicacion $\rightarrow$ HAL de telemetria
$\rightarrow$ USB/UART $\rightarrow$ Processing
\end{center}

Cada flecha representa una interfaz. El modulo de la izquierda ofrece una
funcion; el de la derecha la llama sin necesitar conocer sus detalles internos.
La meta es que los algoritmos que no dependen del hardware se escriban una sola
vez y que cada placa implemente unicamente los accesos fisicos que necesita.

Las palabras ``compartido'' y ``comun'' significan que el codigo se usa al
compilar para ambas placas. No significan que los perifericos internos de ESP32
y EDU-CIAA sean iguales. El MPU6050, el VL53L0X, las unidades y el protocolo de
la aplicacion son los mismos; las bibliotecas de bajo nivel y el arranque no.

%==============================================================================
\section{Que hace el firmware hoy}
%==============================================================================

La version actual realiza un ciclo de percepcion V0:

\begin{enumerate}
  \item Inicializa la IMU MPU6050 y el sensor de distancia VL53L0X.
  \item Lee la velocidad angular del eje Z de la IMU y una distancia del ToF.
  \item Corrige e integra la velocidad angular para estimar el angulo $\theta$.
  \item Acepta distancias entre 0.05 m y 1.50 m; el resto se representa como
        $-1$ (medida invalida o sensor no disponible).
  \item Publica una linea CSV con $\theta$ en radianes y distancia en metros.
  \item Processing representa la orientacion y acumula puntos del ToF mientras
        el robot gira en el lugar.
\end{enumerate}

No hay todavia odometria XY, encoders, motores, navegacion ni grilla de
ocupacion. El angulo proviene de integrar el giroscopio, por lo que deriva con
el tiempo. Esa limitacion es conocida y no equivale a un sistema SLAM.

%==============================================================================
\section{Mapa del workspace}
%==============================================================================

La raiz del repositorio es \ruta{C:/Arturito}. El archivo
\ruta{Arturito.code-workspace} abre cinco carpetas en VS Code:

\begin{lstlisting}
Arturito/
  Arturito.code-workspace
  README.md
  docs/
    ARTURITO MAIN.tex
    CONTEXTO_PROYECTO_ROVER.tex
    % (Diseno preliminar de firmware.tex)
    informe_avance.tex
    explicando_firmware.tex
  shared/
        config/config.h
    layer1_hal/
                  hal_imu.h
                  hal_time.h
                  hal_tof.h
                  hal_telemetry.h
    layer2_drivers/
                  driver_mpu6050.c/.h
                  driver_distancia.c/.h
    layer3_app/
                  app_runtime.c/.h
  esp32/
    platformio.ini
    src/
      main.cpp
      layer1_hal/
                        hal_imu_esp32.cpp
                        hal_time_esp32.cpp
                        hal_tof_esp32.cpp
                        hal_telemetry_esp32.cpp
      tests/
                        imu_only.cpp
                        tof_only.cpp
  edu-ciaa/
            Makefile, board.mk y program.mk
    rover/
      config.mk
      inc/
                        FreeRTOSConfig.h
      src/
        main.c
        layer1_hal/
                              hal_imu_educiaa.cpp
                              hal_time_educiaa.cpp
                              hal_tof_educiaa.cpp
                              hal_telemetry_educiaa.cpp
  processing/
    flecha/
    flecha_y_mapa_v1/
\end{lstlisting}

\begin{description}
  \item[\ruta{docs/}] Documentos del equipo. \ruta{informe_avance.tex} es la
        referencia viva de decisiones y estado; MAIN se conserva como entrega
        academica congelada.
  \item[\ruta{shared/}] Contratos HAL, drivers y logica de aplicacion que no
        deben depender de Arduino ni de sAPI.
  \item[\ruta{esp32/}] Firmware PlatformIO para ESP32, sus HAL de placa y
        programas de diagnostico independientes.
  \item[\ruta{edu-ciaa/}] Proyecto Make/LPCOpen/sAPI para EDU-CIAA. Adapta la
        misma aplicacion y drivers a su arquitectura.
  \item[\ruta{processing/}] Sketches que reciben el CSV por puerto serie y lo
        visualizan.
\end{description}

Las carpetas \ruta{layer3\_robotics/} y \ruta{layer4\_app/} pertenecian a una
organizacion anterior y fueron retiradas. La arquitectura acordada tiene tres
capas; su carpeta de aplicacion es \ruta{shared/layer3\_app/}.

%==============================================================================
\section{Las tres capas}
%==============================================================================

\subsection{Capa 1: HAL de perifericos}

HAL significa \emph{Hardware Abstraction Layer}, o capa de abstraccion de
hardware. Convierte operaciones de una placa concreta en funciones estables
que el resto del programa puede usar.

Los contratos publicos viven en \ruta{shared/layer1\_hal/}. Cada placa tiene
una implementacion con el mismo nombre de funcion:

\begin{itemize}
  \item \texttt{hal\_imu\_init()} inicializa la IMU y devuelve si pudo hacerlo.
  \item \texttt{hal\_imu\_get\_gyro\_z\_rads()} devuelve la velocidad angular
        Z en radianes por segundo.
  \item \texttt{hal\_tof\_init()} inicializa el sensor ToF.
  \item \texttt{hal\_tof\_get\_distance\_m()} devuelve metros, o $-1$ si la
        medida no esta disponible.
  \item \texttt{hal\_millis()} y \texttt{hal\_delay()} ofrecen tiempo y espera
        sin que los drivers llamen directamente a la API de una placa.
  \item \texttt{hal\_telemetry\_write()} envia una linea de texto por el canal
        serie de la placa.
\end{itemize}

La HAL no calcula el angulo del robot, no filtra la distancia y no arma el
mapa. Su responsabilidad es inicializar y leer el hardware, devolviendo
magnitudes con unidades acordadas.

\subsection{Capa 2: drivers}

Los drivers estan en \ruta{shared/layer2\_drivers/}. Conocen el significado
fisico de las muestras, pero usan solamente los contratos HAL.

\begin{itemize}
  \item \texttt{driver\_mpu6050.c} calibra el offset del giroscopio, aplica una
        zona muerta y convierte las muestras en un angulo integrado.
  \item \texttt{driver\_distancia.c} valida que una distancia este dentro del
        rango de trabajo del proyecto.
\end{itemize}

Al estar escritos contra las interfaces compartidas, estos drivers se compilan
sin duplicarlos para ESP32 y EDU-CIAA. La HAL correspondiente queda enlazada
por el build de cada placa.

\subsection{Capa 3: aplicacion compartida}

\ruta{shared/layer3\_app/app\_runtime.c} implementa el runtime FreeRTOS comun:
la inicializacion del sistema y las tareas de adquisicion, actualizacion de
drivers y publicacion. No accede a registros I2C ni llama directamente a
Arduino o sAPI.

La capa 3 todavia no implementa cinematica diferencial, pose XY, grilla ni
navegacion. Esas funciones dependen de encoders y decisiones de V1.

%==============================================================================
\section{Contratos compartidos}
%==============================================================================

\subsection{Configuracion comun}

\ruta{shared/config/config.h} define valores que deben coincidir en ambos
builds:

\begin{itemize}
  \item \texttt{SAMPLE\_PERIOD\_MS = 30}: periodo nominal de las tareas ciclicas.
  \item \texttt{IMU\_I2C\_ADDR = 0x68}: direccion I2C del MPU6050.
  \item \texttt{DIST\_MIN\_M = 0.05} y \texttt{DIST\_MAX\_M = 1.50}: limites
        aceptados para distancia.
\end{itemize}

El build define la placa con \texttt{BOARD\_ESP32} o
\texttt{BOARD\_EDUCIAA}. El header comun no elige hardware por su cuenta.

\subsection{Por que los headers tienen extern C}

Los contratos HAL y drivers estan escritos para poder llamarse desde C y C++.
Los headers rodean sus declaraciones con \texttt{extern "C"} cuando los incluye
C++, evitando que el compilador cambie los nombres de las funciones (\emph{name
mangling}). Asi, el codigo C compartido puede enlazarse con las implementaciones
HAL escritas en C++.

\subsection{Header agrupador}

\ruta{shared/layer1\_hal/hal\_sensors.h} incluye los headers separados de tiempo,
IMU y ToF. Es un header agrupador de compatibilidad; las funciones nuevas
incluyen directamente el contrato especifico que necesitan.

%==============================================================================
\section{Que hace cada modulo}
%==============================================================================

\subsection{Runtime: app\_runtime.c}

\begin{description}
  \item[\texttt{app\_runtime\_start()}] Inicializa el estado en valores seguros
        y crea la tarea bootstrap. Devuelve si FreeRTOS pudo crearla; no ejecuta
        la inicializacion de sensores directamente.
  \item[\texttt{task\_start\_application()}] Publica un encabezado de protocolo,
        inicializa IMU y ToF por separado, calibra el MPU6050, crea mutexes y
        luego crea las tres tareas ciclicas. Si la IMU falla, se detiene el
        runtime de sensores; si falla el ToF, se continua con distancia $-1$.
  \item[\texttt{task\_acquire()}] Lee gyro Z y distancia cruda por HAL y guarda
        una instantanea protegida por mutex. Espera el periodo nominal.
  \item[\texttt{task\_update\_drivers()}] Copia la ultima muestra cruda, integra
        gyro para obtener $\theta$, valida la distancia y guarda el estado de
        aplicacion bajo un segundo mutex.
  \item[\texttt{task\_publish()}] Copia el estado, forma
        \texttt{theta,distancia} con cuatro decimales y lo envia mediante la HAL
        de telemetria. Los retardos son nominales; no constituyen sincronizacion
        exacta de fase.
  \item[\texttt{stack\_depth\_from\_bytes()}] Recibe tamanos conceptuales en
        bytes. En ESP32 conserva ese valor, porque su port de FreeRTOS usa bytes;
        en CIAA lo convierte a palabras \texttt{StackType\_t}, como requiere el
        port estandar ARM.
  \item[\texttt{delete\_task\_if\_created()}] Libera tareas ya creadas si la
        creacion de una tarea posterior falla.
\end{description}

Los mutexes protegen las copias compartidas: el lector no observa una estructura
mientras otra tarea la esta reemplazando. No protegen los buses I2C como un
recurso compartido general; la aplicacion actual asigna las lecturas al flujo
de adquisicion.

\subsection{Driver MPU6050}

\begin{description}
  \item[\texttt{mpu6050\_init()}] Actualmente no hace trabajo: la
        inicializacion fisica ya la hizo \texttt{hal\_imu\_init()}.
  \item[\texttt{mpu6050\_calibrate()}] Espera que el robot quede quieto,
        promedia 1000 lecturas separadas por 2 ms para estimar el offset, y
        reinicia el estado de integracion.
  \item[\texttt{mpu6050\_update\_gyro(gz)}] Resta el offset, elimina velocidades
        dentro de $\pm1^\circ$/s, calcula un promedio trapezoidal entre muestras
        y acumula el angulo con el tiempo transcurrido.
  \item[\texttt{mpu6050\_update()}] Conveniencia para un super-loop o test:
        lee la HAL y delega en \texttt{mpu6050\_update\_gyro()}.
  \item[\texttt{mpu6050\_get\_theta\_rad()}] Devuelve el angulo integrado en
        radianes. La integracion de gyro deriva; no es una pose absoluta.
\end{description}

\subsection{Driver de distancia}

\begin{description}
  \item[\texttt{distancia\_init()}] Actualmente no hace trabajo; la HAL
        inicializa el sensor fisico.
  \item[\texttt{distancia\_validate\_m(d)}] Conserva $d$ si esta entre los
        limites inclusivos; devuelve $-1$ si esta fuera de rango.
  \item[\texttt{distancia\_get\_m()}] Lee la HAL ToF y valida el resultado en
        una sola llamada.
\end{description}

\subsection{Arranques de cada plataforma}

\begin{description}
  \item[ESP32, \texttt{setup()}] Arduino inicia el runtime FreeRTOS antes de
        llamar a \texttt{setup()}. El punto de entrada configura el baud rate y
        solicita \texttt{app\_runtime\_start()}. \texttt{loop()} queda vacio
        porque las tareas ejecutan la aplicacion.
  \item[EDU-CIAA, \texttt{main()}] Ejecuta \texttt{boardConfig()}, crea la tarea
        bootstrap con \texttt{app\_runtime\_start()} y llama a
        \texttt{vTaskStartScheduler()}. La inicializacion y calibracion ocurren
        despues, desde la tarea bootstrap. Si el scheduler retorna, informa un
        error porque normalmente no deberia retornar.
\end{description}

%==============================================================================
\section{FreeRTOS en las dos placas}
%==============================================================================

Las dos placas ejecutan la misma aplicacion FreeRTOS. Las cuatro tareas tienen
responsabilidades separadas:

\begin{longtable}{>{\raggedright\arraybackslash}p{0.27\textwidth}
                  >{\raggedright\arraybackslash}p{0.63\textwidth}}
\textbf{Tarea} & \textbf{Responsabilidad} \\
\hline
\texttt{app\_start} & Inicializacion de sensores, calibracion y creacion de
mutexes/tareas. Se elimina cuando termina el arranque. \\
\texttt{adquirir} & Lee IMU y ToF por HAL y publica una muestra cruda compartida. \\
\texttt{drivers} & Procesa las muestras crudas: offset, zona muerta, integracion
de $\theta$ y validacion de distancia. \\
\texttt{publicar} & Serializa el estado como CSV y delega la escritura a la HAL
de telemetria. \\
\end{longtable}

Las tareas ciclicas usan periodos nominales de 30 ms. Las prioridades son
arranque 4, adquisicion 3, drivers 2 y publicacion 1. CIAA admite hasta siete
niveles en su configuracion actual. Los valores de stack se expresan en bytes en
el runtime y se convierten para el ABI de FreeRTOS de CIAA.

La tarea de adquisicion realiza lecturas I2C secuenciales. La aplicacion no
promete que todas las etapas tengan latencia identica en ambas placas: el ToF
EDU-CIAA usa medicion single-shot y puede ocupar mas tiempo que una lectura
instantanea. FreeRTOS permite que otras tareas progresen, pero no vuelve
instantaneo al periferico.

\subsection{Scheduler y arranque}

\begin{itemize}
  \item En ESP32, Arduino/ESP-IDF arranca FreeRTOS. Crear tareas desde
        \texttt{setup()} no inicia un segundo scheduler.
  \item En EDU-CIAA, \texttt{main()} prepara el hardware, crea el bootstrap y
        llama a \texttt{vTaskStartScheduler()}.
  \item El bootstrap corre con el scheduler activo en ambas placas. Esto es
        necesario porque la calibracion usa \texttt{hal\_delay()}.
  \item CIAA compila el kernel FreeRTOS del repositorio, activa asignacion
        dinamica y utiliza \texttt{heap\_4}. En el estado actual, el heap se
        configura en 24 KiB, las prioridades disponibles son 7 y el tick es de
        1 kHz. La configuracion debe respetar la memoria disponible de LPC4337.
\end{itemize}

Los entry points son distintos porque Arduino y el proyecto LPCOpen/sAPI tienen
ciclos de inicio diferentes. Lo que se comparte es el runtime de aplicacion,
las responsabilidades de las tareas y los contratos; no se fuerza a ambas
plataformas a tener el mismo \texttt{main()} de sistema.

%==============================================================================
\section{HAL de ESP32}
%==============================================================================

\subsection{IMU: hal\_imu\_esp32.cpp}

Usa \texttt{Wire} para I2C. \texttt{hal\_imu\_init()} activa el MPU6050,
consulta su registro de identidad y falla si no recibe una respuesta valida.
\texttt{hal\_imu\_get\_gyro\_z\_rads()} lee los registros de gyro Z, aplica
la sensibilidad de $131$ cuentas por grado/segundo para el rango configurado y
convierte a radianes por segundo.

\subsection{Tiempo: hal\_time\_esp32.cpp}

Adapta \texttt{hal\_delay()} a \texttt{delay()} de Arduino y
\texttt{hal\_millis()} a \texttt{millis()}. El driver MPU usa estas funciones
para calibrar y calcular el intervalo $dt$.

\subsection{ToF: hal\_tof\_esp32.cpp}

Usa la biblioteca Adafruit \texttt{VL53L0X}. \texttt{hal\_tof\_init()} guarda
si la biblioteca encontro e inicializo el sensor. La lectura devuelve metros;
si el dispositivo no inicio o el estado de rango reporta error, devuelve $-1$.

\subsection{Telemetria: hal\_telemetry\_esp32.cpp}

Escribe cada linea con \texttt{Serial.println()}. No genera el contenido del
CSV: recibe texto ya formado desde la aplicacion.

%==============================================================================
\section{HAL de EDU-CIAA}
%==============================================================================

\subsection{IMU y tiempo}

\texttt{hal\_imu\_educiaa.cpp} utiliza las funciones MPU60X0 de sAPI para
inicializar, actualizar y leer gyro Z. \texttt{hal\_time\_educiaa.cpp} adapta
\texttt{hal\_delay()} a la espera sAPI y \texttt{hal\_millis()} a
\texttt{tickRead()}. Con el tick configurado a 1 kHz, el conteo representa
milisegundos.

\subsection{ToF EDU-CIAA}

\texttt{hal\_tof\_educiaa.cpp} implementa el acceso al VL53L0X mediante I2C de
sAPI. A diferencia de ESP32, no usa la biblioteca Adafruit: reproduce la
secuencia de registros del driver de referencia para la medicion single-shot.

Los auxiliares internos leen/escriben registros de 8 y 16 bits, traducen
unidades de timeout, consultan la secuencia de medicion y ejecutan la
configuracion SPAD/tuning/calibracion de referencia. La inicializacion comprueba
identificacion de modelo/revision y limita esperas mediante timeout. La lectura
inicia una medicion, espera su terminacion, lee milimetros y devuelve metros.
Ante fallo de identificacion, calibracion o timeout, marca el ToF no disponible
para que la aplicacion publique $-1$ y la IMU continue.

Esta implementacion compila y enlaza para EDU-CIAA. Sin sensor y placa conectados
no se puede asegurar que la secuencia I2C funcione en el montaje real.

\subsection{Telemetria: hal\_telemetry\_educiaa.cpp}

Implementa el mismo contrato \texttt{hal\_telemetry\_write()} con
\texttt{printf()} de sAPI y terminador CRLF. Asi el runtime compartido no sabe
si debajo escribe Arduino Serial o la consola serie de CIAA.

%==============================================================================
\section{Builds y perfiles}
%==============================================================================

\subsection{ESP32 / PlatformIO}

\ruta{esp32/platformio.ini} establece \texttt{esp32dev} como entorno por
defecto:

\begin{lstlisting}
cd esp32
pio run
pio run -t upload
\end{lstlisting}

Los perfiles \texttt{imu\_only} y \texttt{tof\_only} son diagnosticos
independientes; se seleccionan explicitamente con
\texttt{pio run -e imu\_only} o \texttt{pio run -e tof\_only}. Un comando de
upload con un perfil reemplaza el programa anterior en la placa: no se mezclan
los tres binarios.

\subsection{EDU-CIAA / Make}

\ruta{edu-ciaa/rover/config.mk} selecciona HAL EDU-CIAA, drivers compartidos y
FreeRTOS. La opcion \texttt{USE\_FREERTOS=y} incorpora el kernel; el tipo de
heap selecciona \texttt{heap\_4}. \ruta{edu-ciaa/rover/inc/FreeRTOSConfig.h}
define prioridades, frecuencia de tick, heap y opciones de chequeo.

\begin{lstlisting}
make -C edu-ciaa clean all
\end{lstlisting}

Este comando compila y enlaza; no carga firmware. El target \texttt{clean}
elimina productos de compilacion de \ruta{edu-ciaa/rover/out/}. No borra los
fuentes.

\subsection{Diagnosticos}

\begin{description}
  \item[\texttt{imu\_only}] Compila IMU + tiempo + driver MPU. Publica
        \texttt{theta\_rad,-1} sin depender del ToF ni del runtime completo.
  \item[\texttt{tof\_only}] Compila ToF + driver de distancia. Publica
        \texttt{0,distancia\_m} para comprobar el sensor de forma aislada.
\end{description}

Estos programas validan perifericos concretos, no prueban la planificacion de
las tres tareas del producto.

%==============================================================================
\section{Protocolo serie y Processing}
%==============================================================================

La aplicacion publica a 115200 baudios. Las lineas informativas comienzan con
\texttt{\#}; los sketch las ignoran. Cada muestra de datos tiene exactamente
dos columnas:

\begin{lstlisting}
theta_rad,distancia_m
\end{lstlisting}

$\theta$ esta en radianes; distancia esta en metros. El valor $-1$ indica que
no hay una medida ToF valida. No hay timestamp, posicion XY ni una tercera
columna.

\ruta{processing/flecha/flecha.pde} muestra la orientacion usando theta.
\ruta{processing/flecha\_y\_mapa\_v1/flecha\_y\_mapa\_v1.pde} lee las dos
columnas, descarta logs y datos no numericos, acepta distancias de 0.05 a 1.50 m
y las escala para dibujar una nube polar de hasta 400 puntos. Asume que el rover
gira en torno al origen; no es una grilla ni un mapa en coordenadas mundo.

Los sketches usan \texttt{COM3} como puerto por defecto. Cada PC debe usar el
puerto que realmente asigna al ESP32. Un puerto serie solo puede estar abierto
por una aplicacion a la vez; cerrar el monitor serie de PlatformIO antes de
abrir Processing.

%==============================================================================
\section{Recorrido de una muestra}
%==============================================================================

\begin{enumerate}
  \item El framework de la placa entra en su arranque: Arduino llama a
        \texttt{setup()}, o CIAA ejecuta \texttt{main()}.
  \item El adaptador de placa configura consola/placa y llama a
        \texttt{app\_runtime\_start()}.
  \item CIAA arranca el scheduler. El bootstrap comun inicializa IMU y ToF,
        calibra el MPU y crea los dos mutexes.
  \item Se crea \texttt{adquirir}. Lee $\omega_z$ y distancia por contratos HAL
        y actualiza \texttt{raw\_sample}.
  \item \texttt{drivers} toma una copia protegida, actualiza $\theta$ con la
        muestra IMU y valida distancia. Guarda \texttt{app\_sample}.
  \item \texttt{publicar} copia ese estado, forma el CSV y llama a
        \texttt{hal\_telemetry\_write()}.
  \item La implementacion serie de la placa escribe la linea; Processing la
        interpreta y dibuja la orientacion/medicion.
\end{enumerate}

Si falla la IMU, el bootstrap informa error y no crea las tareas de datos. Si
falla solo el ToF, la tarea de adquisicion sigue leyendo la IMU y el driver de
distancia transforma las medidas no disponibles en $-1$.

%==============================================================================
\section{Estado de validacion y limites}
%==============================================================================

\begin{itemize}
  \item Los entornos PlatformIO \texttt{esp32dev}, \texttt{imu\_only} y
        \texttt{tof\_only} compilan.
  \item El firmware ESP32 fue cargado y el plot V0 habia sido validado con
        Processing antes de compartir el runtime FreeRTOS.
  \item El build limpio EDU-CIAA con FreeRTOS compila y enlaza. No se cargo ni
        se probo con sensores en la placa EDU-CIAA.
  \item La carga del runtime compartido en ESP32 requiere confirmar telemetria
        en serie despues de la ultima modificacion; no se considera prueba de
        ejecucion solo porque haya compilado.
\end{itemize}

Limites funcionales actuales:

\begin{itemize}
  \item El angulo deriva porque se integra un giroscopio; no hay fusion con
        acelerometro ni correccion absoluta.
  \item Sin AS5600 no se calcula desplazamiento XY. El robot se supone en el
        origen para el plot polar V0.
  \item No hay motores, navegacion, grilla V1 ni SLAM implementados.
  \item La inicializacion ToF y los tiempos de medicion de las dos HAL son
        diferentes y requieren validacion fisica en cada placa.
\end{itemize}

%==============================================================================
\section{Glosario}
%==============================================================================

\begin{description}
  \item[HAL] Capa que traduce una API estable a operaciones de hardware de una
        placa concreta.
  \item[Driver] Codigo que transforma muestras HAL en magnitudes utiles para
        la aplicacion, aplicando calibracion, filtrado o validacion.
  \item[FreeRTOS] Sistema operativo en tiempo real que planifica las tareas.
  \item[Scheduler] Planificador que decide que tarea ejecuta el procesador.
  \item[Tarea] Funcion que FreeRTOS ejecuta como flujo independiente, normalmente
        dentro de un bucle.
  \item[Mutex] Cerradura que coordina acceso a datos compartidos entre tareas.
  \item[ToF] Sensor de distancia por tiempo de vuelo; aqui, VL53L0X.
  \item[CSV] Datos separados por comas; el protocolo V0 tiene dos columnas.
  \item[Stack] Memoria temporal de cada tarea para llamadas y variables locales.
\end{description}

% Sección incluida desde ARTURITO MAIN.tex con \documentclass[11pt,a4paper]{article}

\usepackage[utf8]{inputenc}
\usepackage[T1]{fontenc}
\usepackage[spanish,es-noquoting]{babel}
\usepackage[a4paper,margin=2.3cm]{geometry}
\usepackage{array}
\usepackage{longtable}
\usepackage{enumitem}
\usepackage{xcolor}
\usepackage{listings}
\usepackage{hyperref}
\usepackage{parskip}

\hypersetup{
  colorlinks=true,
  linkcolor=black,
  urlcolor=blue!45!black,
  pdftitle={Arturito --- Explicacion del firmware},
  pdfauthor={Equipo Arturito}
}

\urlstyle{tt}
\setlist{itemsep=2pt,topsep=4pt}
\lstset{
  basicstyle=\ttfamily\small,
  columns=fullflexible,
  keepspaces=true,
  breaklines=true,
  frame=single,
  rulecolor=\color{black!25}
}

\newcommand{\ruta}[1]{\path{#1}}
\newcommand{\capa}[1]{\textbf{Capa #1}}

\title{Arturito --- Explicacion del firmware\\
       \large Guia conceptual del workspace, los modulos y el flujo de datos}
\author{Equipo Arturito}
\date{3 de octubre de 2026}

\begin{document}
\maketitle

\begin{abstract}
Este documento explica como esta organizado y como funciona el firmware del
rover Arturito. Esta escrito para alguien que empieza desde cero: presenta el
workspace, las tres capas de software, las diferencias entre ESP32 y EDU-CIAA,
el runtime FreeRTOS compartido, el recorrido de una muestra desde los sensores
hasta Processing y las formas de compilar cada variante.

La explicacion describe el codigo actual. ESP32 ya fue probado con Processing.
La version EDU-CIAA compila y enlaza, pero todavia necesita una prueba fisica de
sensores en la placa. El diseno academico original se conserva por separado;
las decisiones vigentes se registran en \ruta{informe_avance.tex}.
\end{abstract}

\tableofcontents
\newpage

%==============================================================================
\section{Como leer esta guia}
%==============================================================================

El firmware puede entenderse como una cadena de responsabilidades:

\begin{center}
Sensor fisico $\rightarrow$ HAL de la placa $\rightarrow$ driver compartido
$\rightarrow$ tareas de aplicacion $\rightarrow$ HAL de telemetria
$\rightarrow$ USB/UART $\rightarrow$ Processing
\end{center}

Cada flecha representa una interfaz. El modulo de la izquierda ofrece una
funcion; el de la derecha la llama sin necesitar conocer sus detalles internos.
La meta es que los algoritmos que no dependen del hardware se escriban una sola
vez y que cada placa implemente unicamente los accesos fisicos que necesita.

Las palabras ``compartido'' y ``comun'' significan que el codigo se usa al
compilar para ambas placas. No significan que los perifericos internos de ESP32
y EDU-CIAA sean iguales. El MPU6050, el VL53L0X, las unidades y el protocolo de
la aplicacion son los mismos; las bibliotecas de bajo nivel y el arranque no.

%==============================================================================
\section{Que hace el firmware hoy}
%==============================================================================

La version actual realiza un ciclo de percepcion V0:

\begin{enumerate}
  \item Inicializa la IMU MPU6050 y el sensor de distancia VL53L0X.
  \item Lee la velocidad angular del eje Z de la IMU y una distancia del ToF.
  \item Corrige e integra la velocidad angular para estimar el angulo $\theta$.
  \item Acepta distancias entre 0.05 m y 1.50 m; el resto se representa como
        $-1$ (medida invalida o sensor no disponible).
  \item Publica una linea CSV con $\theta$ en radianes y distancia en metros.
  \item Processing representa la orientacion y acumula puntos del ToF mientras
        el robot gira en el lugar.
\end{enumerate}

No hay todavia odometria XY, encoders, motores, navegacion ni grilla de
ocupacion. El angulo proviene de integrar el giroscopio, por lo que deriva con
el tiempo. Esa limitacion es conocida y no equivale a un sistema SLAM.

%==============================================================================
\section{Mapa del workspace}
%==============================================================================

La raiz del repositorio es \ruta{C:/Arturito}. El archivo
\ruta{Arturito.code-workspace} abre cinco carpetas en VS Code:

\begin{lstlisting}
Arturito/
  Arturito.code-workspace
  README.md
  docs/
    ARTURITO MAIN.tex
    CONTEXTO_PROYECTO_ROVER.tex
    % (Diseno preliminar de firmware.tex)
    informe_avance.tex
    explicando_firmware.tex
  shared/
        config/config.h
    layer1_hal/
                  hal_imu.h
                  hal_time.h
                  hal_tof.h
                  hal_telemetry.h
    layer2_drivers/
                  driver_mpu6050.c/.h
                  driver_distancia.c/.h
    layer3_app/
                  app_runtime.c/.h
  esp32/
    platformio.ini
    src/
      main.cpp
      layer1_hal/
                        hal_imu_esp32.cpp
                        hal_time_esp32.cpp
                        hal_tof_esp32.cpp
                        hal_telemetry_esp32.cpp
      tests/
                        imu_only.cpp
                        tof_only.cpp
  edu-ciaa/
            Makefile, board.mk y program.mk
    rover/
      config.mk
      inc/
                        FreeRTOSConfig.h
      src/
        main.c
        layer1_hal/
                              hal_imu_educiaa.cpp
                              hal_time_educiaa.cpp
                              hal_tof_educiaa.cpp
                              hal_telemetry_educiaa.cpp
  processing/
    flecha/
    flecha_y_mapa_v1/
\end{lstlisting}

\begin{description}
  \item[\ruta{docs/}] Documentos del equipo. \ruta{informe_avance.tex} es la
        referencia viva de decisiones y estado; MAIN se conserva como entrega
        academica congelada.
  \item[\ruta{shared/}] Contratos HAL, drivers y logica de aplicacion que no
        deben depender de Arduino ni de sAPI.
  \item[\ruta{esp32/}] Firmware PlatformIO para ESP32, sus HAL de placa y
        programas de diagnostico independientes.
  \item[\ruta{edu-ciaa/}] Proyecto Make/LPCOpen/sAPI para EDU-CIAA. Adapta la
        misma aplicacion y drivers a su arquitectura.
  \item[\ruta{processing/}] Sketches que reciben el CSV por puerto serie y lo
        visualizan.
\end{description}

Las carpetas \ruta{layer3\_robotics/} y \ruta{layer4\_app/} pertenecian a una
organizacion anterior y fueron retiradas. La arquitectura acordada tiene tres
capas; su carpeta de aplicacion es \ruta{shared/layer3\_app/}.

%==============================================================================
\section{Las tres capas}
%==============================================================================

\subsection{Capa 1: HAL de perifericos}

HAL significa \emph{Hardware Abstraction Layer}, o capa de abstraccion de
hardware. Convierte operaciones de una placa concreta en funciones estables
que el resto del programa puede usar.

Los contratos publicos viven en \ruta{shared/layer1\_hal/}. Cada placa tiene
una implementacion con el mismo nombre de funcion:

\begin{itemize}
  \item \texttt{hal\_imu\_init()} inicializa la IMU y devuelve si pudo hacerlo.
  \item \texttt{hal\_imu\_get\_gyro\_z\_rads()} devuelve la velocidad angular
        Z en radianes por segundo.
  \item \texttt{hal\_tof\_init()} inicializa el sensor ToF.
  \item \texttt{hal\_tof\_get\_distance\_m()} devuelve metros, o $-1$ si la
        medida no esta disponible.
  \item \texttt{hal\_millis()} y \texttt{hal\_delay()} ofrecen tiempo y espera
        sin que los drivers llamen directamente a la API de una placa.
  \item \texttt{hal\_telemetry\_write()} envia una linea de texto por el canal
        serie de la placa.
\end{itemize}

La HAL no calcula el angulo del robot, no filtra la distancia y no arma el
mapa. Su responsabilidad es inicializar y leer el hardware, devolviendo
magnitudes con unidades acordadas.

\subsection{Capa 2: drivers}

Los drivers estan en \ruta{shared/layer2\_drivers/}. Conocen el significado
fisico de las muestras, pero usan solamente los contratos HAL.

\begin{itemize}
  \item \texttt{driver\_mpu6050.c} calibra el offset del giroscopio, aplica una
        zona muerta y convierte las muestras en un angulo integrado.
  \item \texttt{driver\_distancia.c} valida que una distancia este dentro del
        rango de trabajo del proyecto.
\end{itemize}

Al estar escritos contra las interfaces compartidas, estos drivers se compilan
sin duplicarlos para ESP32 y EDU-CIAA. La HAL correspondiente queda enlazada
por el build de cada placa.

\subsection{Capa 3: aplicacion compartida}

\ruta{shared/layer3\_app/app\_runtime.c} implementa el runtime FreeRTOS comun:
la inicializacion del sistema y las tareas de adquisicion, actualizacion de
drivers y publicacion. No accede a registros I2C ni llama directamente a
Arduino o sAPI.

La capa 3 todavia no implementa cinematica diferencial, pose XY, grilla ni
navegacion. Esas funciones dependen de encoders y decisiones de V1.

%==============================================================================
\section{Contratos compartidos}
%==============================================================================

\subsection{Configuracion comun}

\ruta{shared/config/config.h} define valores que deben coincidir en ambos
builds:

\begin{itemize}
  \item \texttt{SAMPLE\_PERIOD\_MS = 30}: periodo nominal de las tareas ciclicas.
  \item \texttt{IMU\_I2C\_ADDR = 0x68}: direccion I2C del MPU6050.
  \item \texttt{DIST\_MIN\_M = 0.05} y \texttt{DIST\_MAX\_M = 1.50}: limites
        aceptados para distancia.
\end{itemize}

El build define la placa con \texttt{BOARD\_ESP32} o
\texttt{BOARD\_EDUCIAA}. El header comun no elige hardware por su cuenta.

\subsection{Por que los headers tienen extern C}

Los contratos HAL y drivers estan escritos para poder llamarse desde C y C++.
Los headers rodean sus declaraciones con \texttt{extern "C"} cuando los incluye
C++, evitando que el compilador cambie los nombres de las funciones (\emph{name
mangling}). Asi, el codigo C compartido puede enlazarse con las implementaciones
HAL escritas en C++.

\subsection{Header agrupador}

\ruta{shared/layer1\_hal/hal\_sensors.h} incluye los headers separados de tiempo,
IMU y ToF. Es un header agrupador de compatibilidad; las funciones nuevas
incluyen directamente el contrato especifico que necesitan.

%==============================================================================
\section{Que hace cada modulo}
%==============================================================================

\subsection{Runtime: app\_runtime.c}

\begin{description}
  \item[\texttt{app\_runtime\_start()}] Inicializa el estado en valores seguros
        y crea la tarea bootstrap. Devuelve si FreeRTOS pudo crearla; no ejecuta
        la inicializacion de sensores directamente.
  \item[\texttt{task\_start\_application()}] Publica un encabezado de protocolo,
        inicializa IMU y ToF por separado, calibra el MPU6050, crea mutexes y
        luego crea las tres tareas ciclicas. Si la IMU falla, se detiene el
        runtime de sensores; si falla el ToF, se continua con distancia $-1$.
  \item[\texttt{task\_acquire()}] Lee gyro Z y distancia cruda por HAL y guarda
        una instantanea protegida por mutex. Espera el periodo nominal.
  \item[\texttt{task\_update\_drivers()}] Copia la ultima muestra cruda, integra
        gyro para obtener $\theta$, valida la distancia y guarda el estado de
        aplicacion bajo un segundo mutex.
  \item[\texttt{task\_publish()}] Copia el estado, forma
        \texttt{theta,distancia} con cuatro decimales y lo envia mediante la HAL
        de telemetria. Los retardos son nominales; no constituyen sincronizacion
        exacta de fase.
  \item[\texttt{stack\_depth\_from\_bytes()}] Recibe tamanos conceptuales en
        bytes. En ESP32 conserva ese valor, porque su port de FreeRTOS usa bytes;
        en CIAA lo convierte a palabras \texttt{StackType\_t}, como requiere el
        port estandar ARM.
  \item[\texttt{delete\_task\_if\_created()}] Libera tareas ya creadas si la
        creacion de una tarea posterior falla.
\end{description}

Los mutexes protegen las copias compartidas: el lector no observa una estructura
mientras otra tarea la esta reemplazando. No protegen los buses I2C como un
recurso compartido general; la aplicacion actual asigna las lecturas al flujo
de adquisicion.

\subsection{Driver MPU6050}

\begin{description}
  \item[\texttt{mpu6050\_init()}] Actualmente no hace trabajo: la
        inicializacion fisica ya la hizo \texttt{hal\_imu\_init()}.
  \item[\texttt{mpu6050\_calibrate()}] Espera que el robot quede quieto,
        promedia 1000 lecturas separadas por 2 ms para estimar el offset, y
        reinicia el estado de integracion.
  \item[\texttt{mpu6050\_update\_gyro(gz)}] Resta el offset, elimina velocidades
        dentro de $\pm1^\circ$/s, calcula un promedio trapezoidal entre muestras
        y acumula el angulo con el tiempo transcurrido.
  \item[\texttt{mpu6050\_update()}] Conveniencia para un super-loop o test:
        lee la HAL y delega en \texttt{mpu6050\_update\_gyro()}.
  \item[\texttt{mpu6050\_get\_theta\_rad()}] Devuelve el angulo integrado en
        radianes. La integracion de gyro deriva; no es una pose absoluta.
\end{description}

\subsection{Driver de distancia}

\begin{description}
  \item[\texttt{distancia\_init()}] Actualmente no hace trabajo; la HAL
        inicializa el sensor fisico.
  \item[\texttt{distancia\_validate\_m(d)}] Conserva $d$ si esta entre los
        limites inclusivos; devuelve $-1$ si esta fuera de rango.
  \item[\texttt{distancia\_get\_m()}] Lee la HAL ToF y valida el resultado en
        una sola llamada.
\end{description}

\subsection{Arranques de cada plataforma}

\begin{description}
  \item[ESP32, \texttt{setup()}] Arduino inicia el runtime FreeRTOS antes de
        llamar a \texttt{setup()}. El punto de entrada configura el baud rate y
        solicita \texttt{app\_runtime\_start()}. \texttt{loop()} queda vacio
        porque las tareas ejecutan la aplicacion.
  \item[EDU-CIAA, \texttt{main()}] Ejecuta \texttt{boardConfig()}, crea la tarea
        bootstrap con \texttt{app\_runtime\_start()} y llama a
        \texttt{vTaskStartScheduler()}. La inicializacion y calibracion ocurren
        despues, desde la tarea bootstrap. Si el scheduler retorna, informa un
        error porque normalmente no deberia retornar.
\end{description}

%==============================================================================
\section{FreeRTOS en las dos placas}
%==============================================================================

Las dos placas ejecutan la misma aplicacion FreeRTOS. Las cuatro tareas tienen
responsabilidades separadas:

\begin{longtable}{>{\raggedright\arraybackslash}p{0.27\textwidth}
                  >{\raggedright\arraybackslash}p{0.63\textwidth}}
\textbf{Tarea} & \textbf{Responsabilidad} \\
\hline
\texttt{app\_start} & Inicializacion de sensores, calibracion y creacion de
mutexes/tareas. Se elimina cuando termina el arranque. \\
\texttt{adquirir} & Lee IMU y ToF por HAL y publica una muestra cruda compartida. \\
\texttt{drivers} & Procesa las muestras crudas: offset, zona muerta, integracion
de $\theta$ y validacion de distancia. \\
\texttt{publicar} & Serializa el estado como CSV y delega la escritura a la HAL
de telemetria. \\
\end{longtable}

Las tareas ciclicas usan periodos nominales de 30 ms. Las prioridades son
arranque 4, adquisicion 3, drivers 2 y publicacion 1. CIAA admite hasta siete
niveles en su configuracion actual. Los valores de stack se expresan en bytes en
el runtime y se convierten para el ABI de FreeRTOS de CIAA.

La tarea de adquisicion realiza lecturas I2C secuenciales. La aplicacion no
promete que todas las etapas tengan latencia identica en ambas placas: el ToF
EDU-CIAA usa medicion single-shot y puede ocupar mas tiempo que una lectura
instantanea. FreeRTOS permite que otras tareas progresen, pero no vuelve
instantaneo al periferico.

\subsection{Scheduler y arranque}

\begin{itemize}
  \item En ESP32, Arduino/ESP-IDF arranca FreeRTOS. Crear tareas desde
        \texttt{setup()} no inicia un segundo scheduler.
  \item En EDU-CIAA, \texttt{main()} prepara el hardware, crea el bootstrap y
        llama a \texttt{vTaskStartScheduler()}.
  \item El bootstrap corre con el scheduler activo en ambas placas. Esto es
        necesario porque la calibracion usa \texttt{hal\_delay()}.
  \item CIAA compila el kernel FreeRTOS del repositorio, activa asignacion
        dinamica y utiliza \texttt{heap\_4}. En el estado actual, el heap se
        configura en 24 KiB, las prioridades disponibles son 7 y el tick es de
        1 kHz. La configuracion debe respetar la memoria disponible de LPC4337.
\end{itemize}

Los entry points son distintos porque Arduino y el proyecto LPCOpen/sAPI tienen
ciclos de inicio diferentes. Lo que se comparte es el runtime de aplicacion,
las responsabilidades de las tareas y los contratos; no se fuerza a ambas
plataformas a tener el mismo \texttt{main()} de sistema.

%==============================================================================
\section{HAL de ESP32}
%==============================================================================

\subsection{IMU: hal\_imu\_esp32.cpp}

Usa \texttt{Wire} para I2C. \texttt{hal\_imu\_init()} activa el MPU6050,
consulta su registro de identidad y falla si no recibe una respuesta valida.
\texttt{hal\_imu\_get\_gyro\_z\_rads()} lee los registros de gyro Z, aplica
la sensibilidad de $131$ cuentas por grado/segundo para el rango configurado y
convierte a radianes por segundo.

\subsection{Tiempo: hal\_time\_esp32.cpp}

Adapta \texttt{hal\_delay()} a \texttt{delay()} de Arduino y
\texttt{hal\_millis()} a \texttt{millis()}. El driver MPU usa estas funciones
para calibrar y calcular el intervalo $dt$.

\subsection{ToF: hal\_tof\_esp32.cpp}

Usa la biblioteca Adafruit \texttt{VL53L0X}. \texttt{hal\_tof\_init()} guarda
si la biblioteca encontro e inicializo el sensor. La lectura devuelve metros;
si el dispositivo no inicio o el estado de rango reporta error, devuelve $-1$.

\subsection{Telemetria: hal\_telemetry\_esp32.cpp}

Escribe cada linea con \texttt{Serial.println()}. No genera el contenido del
CSV: recibe texto ya formado desde la aplicacion.

%==============================================================================
\section{HAL de EDU-CIAA}
%==============================================================================

\subsection{IMU y tiempo}

\texttt{hal\_imu\_educiaa.cpp} utiliza las funciones MPU60X0 de sAPI para
inicializar, actualizar y leer gyro Z. \texttt{hal\_time\_educiaa.cpp} adapta
\texttt{hal\_delay()} a la espera sAPI y \texttt{hal\_millis()} a
\texttt{tickRead()}. Con el tick configurado a 1 kHz, el conteo representa
milisegundos.

\subsection{ToF EDU-CIAA}

\texttt{hal\_tof\_educiaa.cpp} implementa el acceso al VL53L0X mediante I2C de
sAPI. A diferencia de ESP32, no usa la biblioteca Adafruit: reproduce la
secuencia de registros del driver de referencia para la medicion single-shot.

Los auxiliares internos leen/escriben registros de 8 y 16 bits, traducen
unidades de timeout, consultan la secuencia de medicion y ejecutan la
configuracion SPAD/tuning/calibracion de referencia. La inicializacion comprueba
identificacion de modelo/revision y limita esperas mediante timeout. La lectura
inicia una medicion, espera su terminacion, lee milimetros y devuelve metros.
Ante fallo de identificacion, calibracion o timeout, marca el ToF no disponible
para que la aplicacion publique $-1$ y la IMU continue.

Esta implementacion compila y enlaza para EDU-CIAA. Sin sensor y placa conectados
no se puede asegurar que la secuencia I2C funcione en el montaje real.

\subsection{Telemetria: hal\_telemetry\_educiaa.cpp}

Implementa el mismo contrato \texttt{hal\_telemetry\_write()} con
\texttt{printf()} de sAPI y terminador CRLF. Asi el runtime compartido no sabe
si debajo escribe Arduino Serial o la consola serie de CIAA.

%==============================================================================
\section{Builds y perfiles}
%==============================================================================

\subsection{ESP32 / PlatformIO}

\ruta{esp32/platformio.ini} establece \texttt{esp32dev} como entorno por
defecto:

\begin{lstlisting}
cd esp32
pio run
pio run -t upload
\end{lstlisting}

Los perfiles \texttt{imu\_only} y \texttt{tof\_only} son diagnosticos
independientes; se seleccionan explicitamente con
\texttt{pio run -e imu\_only} o \texttt{pio run -e tof\_only}. Un comando de
upload con un perfil reemplaza el programa anterior en la placa: no se mezclan
los tres binarios.

\subsection{EDU-CIAA / Make}

\ruta{edu-ciaa/rover/config.mk} selecciona HAL EDU-CIAA, drivers compartidos y
FreeRTOS. La opcion \texttt{USE\_FREERTOS=y} incorpora el kernel; el tipo de
heap selecciona \texttt{heap\_4}. \ruta{edu-ciaa/rover/inc/FreeRTOSConfig.h}
define prioridades, frecuencia de tick, heap y opciones de chequeo.

\begin{lstlisting}
make -C edu-ciaa clean all
\end{lstlisting}

Este comando compila y enlaza; no carga firmware. El target \texttt{clean}
elimina productos de compilacion de \ruta{edu-ciaa/rover/out/}. No borra los
fuentes.

\subsection{Diagnosticos}

\begin{description}
  \item[\texttt{imu\_only}] Compila IMU + tiempo + driver MPU. Publica
        \texttt{theta\_rad,-1} sin depender del ToF ni del runtime completo.
  \item[\texttt{tof\_only}] Compila ToF + driver de distancia. Publica
        \texttt{0,distancia\_m} para comprobar el sensor de forma aislada.
\end{description}

Estos programas validan perifericos concretos, no prueban la planificacion de
las tres tareas del producto.

%==============================================================================
\section{Protocolo serie y Processing}
%==============================================================================

La aplicacion publica a 115200 baudios. Las lineas informativas comienzan con
\texttt{\#}; los sketch las ignoran. Cada muestra de datos tiene exactamente
dos columnas:

\begin{lstlisting}
theta_rad,distancia_m
\end{lstlisting}

$\theta$ esta en radianes; distancia esta en metros. El valor $-1$ indica que
no hay una medida ToF valida. No hay timestamp, posicion XY ni una tercera
columna.

\ruta{processing/flecha/flecha.pde} muestra la orientacion usando theta.
\ruta{processing/flecha\_y\_mapa\_v1/flecha\_y\_mapa\_v1.pde} lee las dos
columnas, descarta logs y datos no numericos, acepta distancias de 0.05 a 1.50 m
y las escala para dibujar una nube polar de hasta 400 puntos. Asume que el rover
gira en torno al origen; no es una grilla ni un mapa en coordenadas mundo.

Los sketches usan \texttt{COM3} como puerto por defecto. Cada PC debe usar el
puerto que realmente asigna al ESP32. Un puerto serie solo puede estar abierto
por una aplicacion a la vez; cerrar el monitor serie de PlatformIO antes de
abrir Processing.

%==============================================================================
\section{Recorrido de una muestra}
%==============================================================================

\begin{enumerate}
  \item El framework de la placa entra en su arranque: Arduino llama a
        \texttt{setup()}, o CIAA ejecuta \texttt{main()}.
  \item El adaptador de placa configura consola/placa y llama a
        \texttt{app\_runtime\_start()}.
  \item CIAA arranca el scheduler. El bootstrap comun inicializa IMU y ToF,
        calibra el MPU y crea los dos mutexes.
  \item Se crea \texttt{adquirir}. Lee $\omega_z$ y distancia por contratos HAL
        y actualiza \texttt{raw\_sample}.
  \item \texttt{drivers} toma una copia protegida, actualiza $\theta$ con la
        muestra IMU y valida distancia. Guarda \texttt{app\_sample}.
  \item \texttt{publicar} copia ese estado, forma el CSV y llama a
        \texttt{hal\_telemetry\_write()}.
  \item La implementacion serie de la placa escribe la linea; Processing la
        interpreta y dibuja la orientacion/medicion.
\end{enumerate}

Si falla la IMU, el bootstrap informa error y no crea las tareas de datos. Si
falla solo el ToF, la tarea de adquisicion sigue leyendo la IMU y el driver de
distancia transforma las medidas no disponibles en $-1$.

%==============================================================================
\section{Estado de validacion y limites}
%==============================================================================

\begin{itemize}
  \item Los entornos PlatformIO \texttt{esp32dev}, \texttt{imu\_only} y
        \texttt{tof\_only} compilan.
  \item El firmware ESP32 fue cargado y el plot V0 habia sido validado con
        Processing antes de compartir el runtime FreeRTOS.
  \item El build limpio EDU-CIAA con FreeRTOS compila y enlaza. No se cargo ni
        se probo con sensores en la placa EDU-CIAA.
  \item La carga del runtime compartido en ESP32 requiere confirmar telemetria
        en serie despues de la ultima modificacion; no se considera prueba de
        ejecucion solo porque haya compilado.
\end{itemize}

Limites funcionales actuales:

\begin{itemize}
  \item El angulo deriva porque se integra un giroscopio; no hay fusion con
        acelerometro ni correccion absoluta.
  \item Sin AS5600 no se calcula desplazamiento XY. El robot se supone en el
        origen para el plot polar V0.
  \item No hay motores, navegacion, grilla V1 ni SLAM implementados.
  \item La inicializacion ToF y los tiempos de medicion de las dos HAL son
        diferentes y requieren validacion fisica en cada placa.
\end{itemize}

%==============================================================================
\section{Glosario}
%==============================================================================

\begin{description}
  \item[HAL] Capa que traduce una API estable a operaciones de hardware de una
        placa concreta.
  \item[Driver] Codigo que transforma muestras HAL en magnitudes utiles para
        la aplicacion, aplicando calibracion, filtrado o validacion.
  \item[FreeRTOS] Sistema operativo en tiempo real que planifica las tareas.
  \item[Scheduler] Planificador que decide que tarea ejecuta el procesador.
  \item[Tarea] Funcion que FreeRTOS ejecuta como flujo independiente, normalmente
        dentro de un bucle.
  \item[Mutex] Cerradura que coordina acceso a datos compartidos entre tareas.
  \item[ToF] Sensor de distancia por tiempo de vuelo; aqui, VL53L0X.
  \item[CSV] Datos separados por comas; el protocolo V0 tiene dos columnas.
  \item[Stack] Memoria temporal de cada tarea para llamadas y variables locales.
\end{description}

% Sección incluida desde ARTURITO MAIN.tex con \input{explicando_firmware.tex}.
% Este archivo contiene un fragmento del informe, no un documento independiente.

% =============================================================================
\section{Diseño del firmware, simulación y depuración}
\label{sec:diseno-firmware}
% =============================================================================

\subsection{Arquitectura propuesta y criterios de diseño}

El firmware se diseñará para integrar la percepción, la estimación de la pose,
la construcción del mapa y el control del rover diferencial. La arquitectura
será portable entre plataformas: el desarrollo y la validación funcional se
realizarán primero sobre ESP32/Arduino ya que no contamos con una EDU-CIAA en nuestras casas para desarrollar y probar software de aplicación; luego se portará a EDU-CIAA reutilizando la
aplicación y los drivers compartidos, y reemplazando la capa de adaptación al
hardware.

Se adopta una organización en tres capas para separar las operaciones propias
de cada placa de los algoritmos del robot. Esta separación permite probar los
módulos por partes, mantener una única implementación de la lógica del rover y
evitar que las decisiones de una placa condicionen el código común. La
arquitectura de ejecución se organizará en tareas periódicas controladas por
tiempo y planificadas por FreeRTOS, la implementación y caracteristicas de este todavia no se decidieron ya que primero optamos por validar los datos del harware e integrarlo con una simple versión del scheduler.

El modelo elegido prioriza interfaces estables entre capas, unidades físicas
comunes y un main compartido. El período de muestreo y las prioridades se
ajustarán mediante mediciones de carga y latencia cuando se integren los
encoders, la construcción del mapa y el control de motores.


\subsection{Módulos organizados en capas}

\begin{description}
  \item[Capa 1 --- HAL de periféricos.] Tendrá una interfaz común y una
        implementación por placa. Inicializará los periféricos y entregará
        magnitudes en unidades acordadas: velocidad angular en rad/s,
        distancia en metros, ángulos de rueda y señales de mando para motores.
        Incluirá HAL para tiempo, MPU6050, VL53L0X, AS5600, GPIO/PWM del DRV8833
        y telemetría. No realizará integración de pose, mapeo ni navegación.
  \item[Capa 2 --- drivers compartidos.] Interpretará las señales físicas de
        cada dispositivo y entregará datos listos para la aplicación. Incluirá lectura, calibración y corrección y validación de datos de los sensores y podrá enviar comandos de generación de señales de
        control para el DRV8833. Sus interfaces no dependerán de Arduino ni de
        SAPI.
  \item[Capa 3 --- aplicación.] Coordinará las tareas FreeRTOS y contendrá la
        lógica propia del rover: adquisición y sincronización de muestras,
        cinemática diferencial y estimación de pose, actualización de la grilla
        de ocupación, selección de consignas según el modo y distribución de
        datos a la navegación y las comunicaciones. Acá viviran los algoritmos de estimación de pose mapa y navegación. 
\end{description}

\subsection{Funcionamiento previsto}
Como esta espcificado 

\begin{enumerate}
  \item Configurar la placa, el scheduler, el canal de telemetría y el estado
        inicial seguro de los actuadores.
  \item Inicializar la IMU MPU6050, el sensor ToF VL53L0X, los encoders AS5600
        y el puente H DRV8833. Informar qué periféricos no están disponibles;
        no habilitar el movimiento si falta una condición necesaria para el
        control seguro.
  \item Calibrar los sensores que lo requieran y poner en cero la pose inicial
        $(x,y,\theta)$. Preparar la representación del mapa correspondiente.
  \item Iniciar las tareas periódicas de FreeRTOS para adquisición,
        procesamiento de la aplicación y publicación/recepción de datos.
  \item En cada ciclo de adquisición, leer gyro Z, distancia ToF y ángulos de
        ambas ruedas; asociar las muestras con el instante de lectura y
        entregarlas a la aplicación.
  \item Procesar las muestras: corregir el gyro, estimar el desplazamiento de
        cada rueda, actualizar la pose del rover y validar la distancia.
  \item Actualizar el mapa con el rayo medido por el ToF y la pose vigente.
 Se actualizará una grilla de ocupación 2D en
        coordenadas $(x,y,\theta)$.
  \item Según el modo seleccionado, obtener la consigna desde navegación
        autónoma o desde teleoperación, convertirla en velocidades de rueda y
        comandar el DRV8833 mediante la HAL. Ante una condición de parada o una
        pérdida de comunicación que supere el límite definido, detener los
        motores.
  \item Publicar pose, estado y datos del mapa para la interfaz externa;
        repetir el ciclo mientras el sistema esté activo.
\end{enumerate}

La navegación autónoma, el control de motores, la integración de los encoders AS5600 y la
grilla de ocupacion forman parte del desarrollo previsto; no se presentan como
funcionalidades ya verificadas. SLAM permanece como objetivo secundario: la
primera grilla se construirá con la pose estimada por odometría.

\subsection{Bibliotecas y herramientas externas previstas}

\begin{itemize}
  \item \textbf{ESP32:} framework Arduino y su integración con FreeRTOS;
        \texttt{Wire} para I2C; biblioteca Adafruit \texttt{VL53L0X} para el
        sensor ToF. Para el transporte inalámbrico se utilizará la capacidad
        disponible en ESP32 una vez elegido el protocolo de enlace.
  \item \textbf{EDU-CIAA:} LPCOpen y SAPI para la placa y sus periféricos, más
        el kernel FreeRTOS configurado en el proyecto. La HAL de ToF implementa
        el acceso I2C por registros y no utiliza la biblioteca Adafruit.
  \item \textbf{Código del proyecto:} contratos HAL, drivers de MPU6050,
        VL53L0X, AS5600 y DRV8833, y módulos de pose, mapa y navegación. Los
        drivers de AS5600 y DRV8833 se incorporarán al contar con el hardware y
        sus interfaces acordadas.
  \item \textbf{PC:} Processing y su biblioteca Serial se utilizan en V0 para
        visualizar las muestras recibidas. La interfaz objetivo mostrará la
        pose del rover y el mapa actualizado.
\end{itemize}

\subsection{Software de PC e interfaz de usuario prevista}

El alcance confirmado para el programa de PC es visualizar la pose del rover y
el mapa actualizado. El enlace serie USB/UART se utiliza en V0; para transmitir
pose y mapa sin cable se prevé incorporar Wi-Fi o Bluetooth en una etapa
posterior. La elección entre ambos protocolos queda pendiente. Los
requerimientos generales también mencionan órdenes START/STOP y, como objetivos
secundarios, teleoperación y planificación de rutas; su incorporación como
controles de esta interfaz queda por definir.

\begin{center}
\fbox{\parbox{0.91\linewidth}{\centering
Inicio $\rightarrow$ conexión con el rover $\rightarrow$ espera de datos
$\rightarrow$ visualización continua de pose y mapa $\rightarrow$ pérdida de
enlace: indicar desconexión / conexión restablecida: reanudar visualización}}
\end{center}

Si se pierde el enlace durante la operación, el firmware deberá aplicar la
parada segura y la interfaz indicará la desconexión. El flujo representa la
interfaz prevista; las pantallas de Processing disponibles actualmente solo
reciben y dibujan telemetría.

\subsection{Ensayos realizados y estado de validación}

Los ensayos se realizaron con el firmware V0 y hardware conectado, usando
programas de prueba y Processing; no se dispone de una simulación integral del
rover. Se separaron perfiles de prueba para compilar la lectura de IMU y la del
ToF de manera independiente. También se implementó \texttt{imu\_compare}, que
transmite gyro Z crudo, ángulo crudo integrado y ángulo procesado para comparar
la integración sin corrección con la que aplica calibración y zona muerta.

El firmware V0 de ESP32 fue cargado y validado con Processing: se comprobó la
recepción de las muestras serie y la visualización de orientación y distancia.
El contrato de aplicación V0 transmite a 115200 baudios dos columnas,
\texttt{theta\_rad,distancia\_m}; una distancia inválida se representa con
$-1$. El build actual con FreeRTOS compartido compila, pero su telemetría debe
confirmarse después de la última modificación. El firmware EDU-CIAA compila y
enlaza, aunque falta probarlo con sensores conectados a la placa.

Los archivos disponibles no contienen capturas de pantalla ni resultados
numéricos de precisión o error. Para documentar visualmente los ensayos se debe
incorporar una captura de Processing durante la prueba V0 y, cuando se realice,
otra del gráfico de comparación de la IMU. Estas evidencias no se reemplazan
por imágenes ilustrativas.




\end{document}.
% Este archivo contiene un fragmento del informe, no un documento independiente.

% =============================================================================
\section{Diseño del firmware, simulación y depuración}
\label{sec:diseno-firmware}
% =============================================================================

\subsection{Arquitectura propuesta y criterios de diseño}

El firmware se diseñará para integrar la percepción, la estimación de la pose,
la construcción del mapa y el control del rover diferencial. La arquitectura
será portable entre plataformas: el desarrollo y la validación funcional se
realizarán primero sobre ESP32/Arduino ya que no contamos con una EDU-CIAA en nuestras casas para desarrollar y probar software de aplicación; luego se portará a EDU-CIAA reutilizando la
aplicación y los drivers compartidos, y reemplazando la capa de adaptación al
hardware.

Se adopta una organización en tres capas para separar las operaciones propias
de cada placa de los algoritmos del robot. Esta separación permite probar los
módulos por partes, mantener una única implementación de la lógica del rover y
evitar que las decisiones de una placa condicionen el código común. La
arquitectura de ejecución se organizará en tareas periódicas controladas por
tiempo y planificadas por FreeRTOS, la implementación y caracteristicas de este todavia no se decidieron ya que primero optamos por validar los datos del harware e integrarlo con una simple versión del scheduler.

El modelo elegido prioriza interfaces estables entre capas, unidades físicas
comunes y un main compartido. El período de muestreo y las prioridades se
ajustarán mediante mediciones de carga y latencia cuando se integren los
encoders, la construcción del mapa y el control de motores.


\subsection{Módulos organizados en capas}

\begin{description}
  \item[Capa 1 --- HAL de periféricos.] Tendrá una interfaz común y una
        implementación por placa. Inicializará los periféricos y entregará
        magnitudes en unidades acordadas: velocidad angular en rad/s,
        distancia en metros, ángulos de rueda y señales de mando para motores.
        Incluirá HAL para tiempo, MPU6050, VL53L0X, AS5600, GPIO/PWM del DRV8833
        y telemetría. No realizará integración de pose, mapeo ni navegación.
  \item[Capa 2 --- drivers compartidos.] Interpretará las señales físicas de
        cada dispositivo y entregará datos listos para la aplicación. Incluirá lectura, calibración y corrección y validación de datos de los sensores y podrá enviar comandos de generación de señales de
        control para el DRV8833. Sus interfaces no dependerán de Arduino ni de
        SAPI.
  \item[Capa 3 --- aplicación.] Coordinará las tareas FreeRTOS y contendrá la
        lógica propia del rover: adquisición y sincronización de muestras,
        cinemática diferencial y estimación de pose, actualización de la grilla
        de ocupación, selección de consignas según el modo y distribución de
        datos a la navegación y las comunicaciones. Acá viviran los algoritmos de estimación de pose mapa y navegación. 
\end{description}

\subsection{Funcionamiento previsto}
Como esta espcificado 

\begin{enumerate}
  \item Configurar la placa, el scheduler, el canal de telemetría y el estado
        inicial seguro de los actuadores.
  \item Inicializar la IMU MPU6050, el sensor ToF VL53L0X, los encoders AS5600
        y el puente H DRV8833. Informar qué periféricos no están disponibles;
        no habilitar el movimiento si falta una condición necesaria para el
        control seguro.
  \item Calibrar los sensores que lo requieran y poner en cero la pose inicial
        $(x,y,\theta)$. Preparar la representación del mapa correspondiente.
  \item Iniciar las tareas periódicas de FreeRTOS para adquisición,
        procesamiento de la aplicación y publicación/recepción de datos.
  \item En cada ciclo de adquisición, leer gyro Z, distancia ToF y ángulos de
        ambas ruedas; asociar las muestras con el instante de lectura y
        entregarlas a la aplicación.
  \item Procesar las muestras: corregir el gyro, estimar el desplazamiento de
        cada rueda, actualizar la pose del rover y validar la distancia.
  \item Actualizar el mapa con el rayo medido por el ToF y la pose vigente.
 Se actualizará una grilla de ocupación 2D en
        coordenadas $(x,y,\theta)$.
  \item Según el modo seleccionado, obtener la consigna desde navegación
        autónoma o desde teleoperación, convertirla en velocidades de rueda y
        comandar el DRV8833 mediante la HAL. Ante una condición de parada o una
        pérdida de comunicación que supere el límite definido, detener los
        motores.
  \item Publicar pose, estado y datos del mapa para la interfaz externa;
        repetir el ciclo mientras el sistema esté activo.
\end{enumerate}

La navegación autónoma, el control de motores, la integración de los encoders AS5600 y la
grilla de ocupacion forman parte del desarrollo previsto; no se presentan como
funcionalidades ya verificadas. SLAM permanece como objetivo secundario: la
primera grilla se construirá con la pose estimada por odometría.

\subsection{Bibliotecas y herramientas externas previstas}

\begin{itemize}
  \item \textbf{ESP32:} framework Arduino y su integración con FreeRTOS;
        \texttt{Wire} para I2C; biblioteca Adafruit \texttt{VL53L0X} para el
        sensor ToF. Para el transporte inalámbrico se utilizará la capacidad
        disponible en ESP32 una vez elegido el protocolo de enlace.
  \item \textbf{EDU-CIAA:} LPCOpen y SAPI para la placa y sus periféricos, más
        el kernel FreeRTOS configurado en el proyecto. La HAL de ToF implementa
        el acceso I2C por registros y no utiliza la biblioteca Adafruit.
  \item \textbf{Código del proyecto:} contratos HAL, drivers de MPU6050,
        VL53L0X, AS5600 y DRV8833, y módulos de pose, mapa y navegación. Los
        drivers de AS5600 y DRV8833 se incorporarán al contar con el hardware y
        sus interfaces acordadas.
  \item \textbf{PC:} Processing y su biblioteca Serial se utilizan en V0 para
        visualizar las muestras recibidas. La interfaz objetivo mostrará la
        pose del rover y el mapa actualizado.
\end{itemize}

\subsection{Software de PC e interfaz de usuario prevista}

El alcance confirmado para el programa de PC es visualizar la pose del rover y
el mapa actualizado. El enlace serie USB/UART se utiliza en V0; para transmitir
pose y mapa sin cable se prevé incorporar Wi-Fi o Bluetooth en una etapa
posterior. La elección entre ambos protocolos queda pendiente. Los
requerimientos generales también mencionan órdenes START/STOP y, como objetivos
secundarios, teleoperación y planificación de rutas; su incorporación como
controles de esta interfaz queda por definir.

\begin{center}
\fbox{\parbox{0.91\linewidth}{\centering
Inicio $\rightarrow$ conexión con el rover $\rightarrow$ espera de datos
$\rightarrow$ visualización continua de pose y mapa $\rightarrow$ pérdida de
enlace: indicar desconexión / conexión restablecida: reanudar visualización}}
\end{center}

Si se pierde el enlace durante la operación, el firmware deberá aplicar la
parada segura y la interfaz indicará la desconexión. El flujo representa la
interfaz prevista; las pantallas de Processing disponibles actualmente solo
reciben y dibujan telemetría.

\subsection{Ensayos realizados y estado de validación}

Los ensayos se realizaron con el firmware V0 y hardware conectado, usando
programas de prueba y Processing; no se dispone de una simulación integral del
rover. Se separaron perfiles de prueba para compilar la lectura de IMU y la del
ToF de manera independiente. También se implementó \texttt{imu\_compare}, que
transmite gyro Z crudo, ángulo crudo integrado y ángulo procesado para comparar
la integración sin corrección con la que aplica calibración y zona muerta.

El firmware V0 de ESP32 fue cargado y validado con Processing: se comprobó la
recepción de las muestras serie y la visualización de orientación y distancia.
El contrato de aplicación V0 transmite a 115200 baudios dos columnas,
\texttt{theta\_rad,distancia\_m}; una distancia inválida se representa con
$-1$. El build actual con FreeRTOS compartido compila, pero su telemetría debe
confirmarse después de la última modificación. El firmware EDU-CIAA compila y
enlaza, aunque falta probarlo con sensores conectados a la placa.

Los archivos disponibles no contienen capturas de pantalla ni resultados
numéricos de precisión o error. Para documentar visualmente los ensayos se debe
incorporar una captura de Processing durante la prueba V0 y, cuando se realice,
otra del gráfico de comparación de la IMU. Estas evidencias no se reemplazan
por imágenes ilustrativas.




\end{document}.
% Este archivo contiene un fragmento del informe, no un documento independiente.

% =============================================================================
\section{Diseño del firmware, simulación y depuración}
\label{sec:diseno-firmware}
% =============================================================================

\subsection{Arquitectura propuesta y criterios de diseño}

El firmware se diseñará para integrar la percepción, la estimación de la pose,
la construcción del mapa y el control del rover diferencial. La arquitectura
será portable entre plataformas: el desarrollo y la validación funcional se
realizarán primero sobre ESP32/Arduino ya que no contamos con una EDU-CIAA en nuestras casas para desarrollar y probar software de aplicación; luego se portará a EDU-CIAA reutilizando la
aplicación y los drivers compartidos, y reemplazando la capa de adaptación al
hardware.

Se adopta una organización en tres capas para separar las operaciones propias
de cada placa de los algoritmos del robot. Esta separación permite probar los
módulos por partes, mantener una única implementación de la lógica del rover y
evitar que las decisiones de una placa condicionen el código común. La
arquitectura de ejecución se organizará en tareas periódicas controladas por
tiempo y planificadas por FreeRTOS, la implementación y caracteristicas de este todavia no se decidieron ya que primero optamos por validar los datos del harware e integrarlo con una simple versión del scheduler.

El modelo elegido prioriza interfaces estables entre capas, unidades físicas
comunes y un main compartido. El período de muestreo y las prioridades se
ajustarán mediante mediciones de carga y latencia cuando se integren los
encoders, la construcción del mapa y el control de motores.


\subsection{Módulos organizados en capas}

\begin{description}
  \item[Capa 1 --- HAL de periféricos.] Tendrá una interfaz común y una
        implementación por placa. Inicializará los periféricos y entregará
        magnitudes en unidades acordadas: velocidad angular en rad/s,
        distancia en metros, ángulos de rueda y señales de mando para motores.
        Incluirá HAL para tiempo, MPU6050, VL53L0X, AS5600, GPIO/PWM del DRV8833
        y telemetría. No realizará integración de pose, mapeo ni navegación.
  \item[Capa 2 --- drivers compartidos.] Interpretará las señales físicas de
        cada dispositivo y entregará datos listos para la aplicación. Incluirá lectura, calibración y corrección y validación de datos de los sensores y podrá enviar comandos de generación de señales de
        control para el DRV8833. Sus interfaces no dependerán de Arduino ni de
        SAPI.
  \item[Capa 3 --- aplicación.] Coordinará las tareas FreeRTOS y contendrá la
        lógica propia del rover: adquisición y sincronización de muestras,
        cinemática diferencial y estimación de pose, actualización de la grilla
        de ocupación, selección de consignas según el modo y distribución de
        datos a la navegación y las comunicaciones. Acá viviran los algoritmos de estimación de pose mapa y navegación. 
\end{description}

\subsection{Funcionamiento previsto}
Como esta espcificado 

\begin{enumerate}
  \item Configurar la placa, el scheduler, el canal de telemetría y el estado
        inicial seguro de los actuadores.
  \item Inicializar la IMU MPU6050, el sensor ToF VL53L0X, los encoders AS5600
        y el puente H DRV8833. Informar qué periféricos no están disponibles;
        no habilitar el movimiento si falta una condición necesaria para el
        control seguro.
  \item Calibrar los sensores que lo requieran y poner en cero la pose inicial
        $(x,y,\theta)$. Preparar la representación del mapa correspondiente.
  \item Iniciar las tareas periódicas de FreeRTOS para adquisición,
        procesamiento de la aplicación y publicación/recepción de datos.
  \item En cada ciclo de adquisición, leer gyro Z, distancia ToF y ángulos de
        ambas ruedas; asociar las muestras con el instante de lectura y
        entregarlas a la aplicación.
  \item Procesar las muestras: corregir el gyro, estimar el desplazamiento de
        cada rueda, actualizar la pose del rover y validar la distancia.
  \item Actualizar el mapa con el rayo medido por el ToF y la pose vigente.
 Se actualizará una grilla de ocupación 2D en
        coordenadas $(x,y,\theta)$.
  \item Según el modo seleccionado, obtener la consigna desde navegación
        autónoma o desde teleoperación, convertirla en velocidades de rueda y
        comandar el DRV8833 mediante la HAL. Ante una condición de parada o una
        pérdida de comunicación que supere el límite definido, detener los
        motores.
  \item Publicar pose, estado y datos del mapa para la interfaz externa;
        repetir el ciclo mientras el sistema esté activo.
\end{enumerate}

La navegación autónoma, el control de motores, la integración de los encoders AS5600 y la
grilla de ocupacion forman parte del desarrollo previsto; no se presentan como
funcionalidades ya verificadas. SLAM permanece como objetivo secundario: la
primera grilla se construirá con la pose estimada por odometría.

\subsection{Bibliotecas y herramientas externas previstas}

\begin{itemize}
  \item \textbf{ESP32:} framework Arduino y su integración con FreeRTOS;
        \texttt{Wire} para I2C; biblioteca Adafruit \texttt{VL53L0X} para el
        sensor ToF. Para el transporte inalámbrico se utilizará la capacidad
        disponible en ESP32 una vez elegido el protocolo de enlace.
  \item \textbf{EDU-CIAA:} LPCOpen y SAPI para la placa y sus periféricos, más
        el kernel FreeRTOS configurado en el proyecto. La HAL de ToF implementa
        el acceso I2C por registros y no utiliza la biblioteca Adafruit.
  \item \textbf{Código del proyecto:} contratos HAL, drivers de MPU6050,
        VL53L0X, AS5600 y DRV8833, y módulos de pose, mapa y navegación. Los
        drivers de AS5600 y DRV8833 se incorporarán al contar con el hardware y
        sus interfaces acordadas.
  \item \textbf{PC:} Processing y su biblioteca Serial se utilizan en V0 para
        visualizar las muestras recibidas. La interfaz objetivo mostrará la
        pose del rover y el mapa actualizado.
\end{itemize}

\subsection{Software de PC e interfaz de usuario prevista}

El alcance confirmado para el programa de PC es visualizar la pose del rover y
el mapa actualizado. El enlace serie USB/UART se utiliza en V0; para transmitir
pose y mapa sin cable se prevé incorporar Wi-Fi o Bluetooth en una etapa
posterior. La elección entre ambos protocolos queda pendiente. Los
requerimientos generales también mencionan órdenes START/STOP y, como objetivos
secundarios, teleoperación y planificación de rutas; su incorporación como
controles de esta interfaz queda por definir.

\begin{center}
\fbox{\parbox{0.91\linewidth}{\centering
Inicio $\rightarrow$ conexión con el rover $\rightarrow$ espera de datos
$\rightarrow$ visualización continua de pose y mapa $\rightarrow$ pérdida de
enlace: indicar desconexión / conexión restablecida: reanudar visualización}}
\end{center}

Si se pierde el enlace durante la operación, el firmware deberá aplicar la
parada segura y la interfaz indicará la desconexión. El flujo representa la
interfaz prevista; las pantallas de Processing disponibles actualmente solo
reciben y dibujan telemetría.

\subsection{Ensayos realizados y estado de validación}

Los ensayos se realizaron con el firmware V0 y hardware conectado, usando
programas de prueba y Processing; no se dispone de una simulación integral del
rover. Se separaron perfiles de prueba para compilar la lectura de IMU y la del
ToF de manera independiente. También se implementó \texttt{imu\_compare}, que
transmite gyro Z crudo, ángulo crudo integrado y ángulo procesado para comparar
la integración sin corrección con la que aplica calibración y zona muerta.

El firmware V0 de ESP32 fue cargado y validado con Processing: se comprobó la
recepción de las muestras serie y la visualización de orientación y distancia.
El contrato de aplicación V0 transmite a 115200 baudios dos columnas,
\texttt{theta\_rad,distancia\_m}; una distancia inválida se representa con
$-1$. El build actual con FreeRTOS compartido compila, pero su telemetría debe
confirmarse después de la última modificación. El firmware EDU-CIAA compila y
enlaza, aunque falta probarlo con sensores conectados a la placa.

Los archivos disponibles no contienen capturas de pantalla ni resultados
numéricos de precisión o error. Para documentar visualmente los ensayos se debe
incorporar una captura de Processing durante la prueba V0 y, cuando se realice,
otra del gráfico de comparación de la IMU. Estas evidencias no se reemplazan
por imágenes ilustrativas.




\end{document}.
% Mantener este archivo como fragmento: no agregar preámbulo ni \begin{document}.

% =============================================================================
\section{Diseño del firmware, simulación y depuración}
\label{sec:diseno-firmware}
% =============================================================================

\subsection{Alcance y arquitectura elegida}

El firmware actual corresponde a una primera versión de percepción (V0). En
esta etapa adquiere la velocidad angular del eje $Z$ de una IMU MPU6050 y una
distancia del sensor ToF VL53L0X, procesa ambas mediciones y transmite el estado
por el puerto serie. El ángulo se obtiene integrando el giroscopio; por ello
puede acumular deriva y no representa una pose absoluta.

La arquitectura es modular, organizada en tres capas, y utiliza tareas
periódicas planificadas por FreeRTOS. No es una arquitectura disparada por
eventos: las tareas repiten su ciclo con un período nominal de
\texttt{SAMPLE\_PERIOD\_MS = 30 ms}. FreeRTOS tiene planificación apropiativa por
prioridades; cada tarea cede el procesador al bloquearse durante la espera. La
adquisición de sensores y el procesamiento no quedan sincronizados en fase, por
lo que los 30 ms son un período nominal y no una garantía de latencia exacta.

\begin{center}
\fbox{\parbox{0.94\linewidth}{\centering
Sensores (I2C) $\rightarrow$ HAL de la placa $\rightarrow$ drivers compartidos
$\rightarrow$ tareas FreeRTOS $\rightarrow$ HAL serie $\rightarrow$ USB/UART
$\rightarrow$ Processing}}
\end{center}

El ciclo de operación puede resumirse con el siguiente pseudocódigo:

\begin{enumerate}
  \item Iniciar la consola serie y crear la tarea de arranque.
  \item Inicializar la MPU6050 y el VL53L0X mediante sus respectivas HAL.
  \item Si la IMU no responde, informar el error y no iniciar la adquisición.
        Si el ToF no responde, continuar e indicar distancia inválida.
  \item Calibrar el offset del giroscopio con el robot quieto y crear los
        mutexes que protegen las muestras compartidas.
  \item Crear las tareas periódicas de adquisición, procesamiento y publicación.
  \item Repetir: leer gyro Z y distancia; guardar una instantánea protegida;
        corregir e integrar gyro para estimar $\theta$; validar distancia;
        transmitir $\theta$ y distancia como una línea CSV.
\end{enumerate}

\subsection{Capas y módulos del firmware}

\begin{description}
  \item[Capa 1 --- HAL de hardware.] Los contratos comunes están en
        \texttt{shared/layer1\_hal/}. Las implementaciones específicas están en
        \texttt{esp32/src/layer1\_hal/} y
        \texttt{edu-ciaa/rover/src/layer1\_hal/}. La HAL de IMU inicializa y lee
        gyro Z en rad/s; la HAL ToF inicializa y devuelve metros o $-1$ si no hay
        una medida disponible; las HAL de tiempo y telemetría adaptan el reloj,
        las esperas y la salida serie de cada placa.
  \item[Capa 2 --- drivers.] Los módulos comunes
        \texttt{shared/layer2\_drivers/driver\_mpu6050.c} y
        \texttt{driver\_distancia.c} no acceden directamente a los periféricos.
        El primero promedia 1000 muestras para estimar el offset, aplica una
        zona muerta de $1^\circ$/s e integra gyro Z mediante integración
        trapezoidal. El segundo acepta distancias entre 0,05 m y 1,50 m,
        inclusive; fuera de ese intervalo informa $-1$.
  \item[Capa 3 --- aplicación.] \texttt{shared/layer3\_app/app\_runtime.c}
        inicializa el sistema, gestiona el estado y crea las tareas. La tarea
        \texttt{adquirir} obtiene las muestras de los sensores; \texttt{drivers}
        procesa la última muestra; y \texttt{publicar} genera la trama serie.
        Dos mutexes protegen las copias de datos crudos y procesados. La tarea
        de arranque inicializa sensores, calibra la IMU y crea las tareas
        cíclicas.
\end{description}

Las prioridades configuradas son 4 para arranque, 3 para adquisición, 2 para
procesamiento y 1 para publicación. Al terminar la inicialización, la tarea de
arranque se elimina. En ESP32, Arduino ya inició FreeRTOS antes de ejecutar
\texttt{setup()}; en EDU-CIAA, \texttt{main()} crea la tarea de arranque y luego
llama al scheduler.

\subsection{Bibliotecas y herramientas externas}

\begin{itemize}
  \item \textbf{ESP32:} framework Arduino para PlatformIO; FreeRTOS provisto
        por el entorno ESP32; \texttt{Wire} para el bus I2C; y biblioteca
        \texttt{Adafruit\_VL53L0X} para el ToF.
  \item \textbf{EDU-CIAA:} LPCOpen y SAPI para la placa y los periféricos,
        junto con el kernel FreeRTOS configurado en el proyecto. La HAL del
        VL53L0X realiza el acceso por registros I2C, sin depender de la
        biblioteca de Adafruit.
  \item \textbf{Código del proyecto:} contratos HAL, drivers MPU6050 y ToF, y
        runtime de aplicación compartidos entre plataformas. No se utiliza
        actualmente un RTOS del proyecto CIAA distinto de FreeRTOS.
  \item \textbf{PC:} Processing y su biblioteca Serial reciben la telemetría
        por USB/UART y dibujan la orientación y las lecturas de distancia.
\end{itemize}

\subsection{Formato de comunicación con el programa de PC}

El firmware de aplicación transmite a 115200 baudios una línea por muestra con
dos campos numéricos separados por coma:

\begin{center}
\texttt{theta\_rad,distancia\_m}
\end{center}

El ángulo se expresa en radianes y la distancia en metros. Las líneas de estado
y error comienzan con \texttt{\#}; el sketch de visualización las descarta. Una
distancia inválida se representa con $-1$. Este formato es un protocolo serie
de prueba V0, sin número de secuencia ni marca de tiempo.

\subsection{Interfaz de usuario y flujo de estados}

La interfaz de usuario actual es el sketch
\texttt{processing/flecha\_y\_mapa\_v1/flecha\_y\_mapa\_v1.pde}. No envía
órdenes al robot: presenta la orientación estimada, la distancia medida y una
nube polar acumulada de hasta 400 puntos. Los puntos se calculan respecto de un
origen fijo, suponiendo que el robot gira en el lugar. Por lo tanto, esta
interfaz permite observar sensores y telemetría, pero no muestra una grilla de
ocupación ni un mapa métrico con pose XY.

\begin{center}
\fbox{\parbox{0.88\linewidth}{\centering
Inicio del sketch $\rightarrow$ selección/configuración del puerto serie
$\rightarrow$ espera de tramas $\rightarrow$ \{línea de estado: actualizar
estado; trama CSV válida: dibujar orientación y distancia; trama inválida:
descartar\} $\rightarrow$ continuar recibiendo}}
\end{center}

El control START/STOP, la teleoperación y la selección entre navegación manual
y autónoma pertenecen a objetivos posteriores. No forman parte de la interfaz
ni del firmware V0 implementado.

\subsection{Ensayos, simulación y depuración}

La verificación realizada hasta ahora es experimental con el hardware y por
consola serie; no se cuenta con una simulación eléctrica o de movimiento del
rover. Para aislar periféricos existen entornos de prueba PlatformIO
\texttt{imu\_only}, \texttt{tof\_only} e \texttt{imu\_compare}. Este último
transmite tres campos (velocidad angular cruda, ángulo crudo integrado y ángulo
procesado) y se utiliza con
\texttt{processing/imu\_compare/imu\_compare.pde} para visualizar la diferencia
entre ambas integraciones. Ese formato de tres columnas es exclusivo de la
prueba y no coincide con el CSV de dos columnas de la aplicación V0.

El sketch \texttt{flecha\_y\_mapa\_v1} consume el CSV de la aplicación normal.
El monitor serie de PlatformIO debe cerrarse antes de abrir Processing, porque
ambos programas no pueden utilizar el mismo puerto simultáneamente. Los mensajes
de arranque permiten distinguir la inicialización de la IMU, el ToF, la
calibración y los errores.

Según el registro de avance del proyecto, los entornos de ESP32
\texttt{esp32dev}, \texttt{imu\_only} y \texttt{tof\_only} compilan, y una
versión anterior del firmware ESP32 fue probada con Processing. La versión
actualizada con runtime FreeRTOS compartido requiere confirmar nuevamente la
telemetría en serie después de su última modificación. El firmware EDU-CIAA
compila y enlaza, pero todavía no fue cargado ni ensayado con sensores en la
placa. No se incluyen capturas de resultados porque no hay imágenes de ensayos
registradas en los archivos del informe al momento de esta redacción.

\subsection{Alcance pendiente}

El firmware actual no implementa encoders, odometría de traslación, control de
motores, navegación autónoma, grilla de ocupación ni SLAM. El sensor ToF entrega
una medición puntual; el barrido angular mecánico y la asociación de cada medida
con una pose global deberán incorporarse para construir el mapa previsto. La
interfaz de usuario deberá ampliarse cuando se implementen los comandos remotos
y la visualización de dicho mapa.
